%%==========================================================================
%% This is the very first version of a Streetwise LaTeX template. By
%% chosing the document class, it is easily convertible to a two-column IEEE
%% paper format, provided any conflicting packages and (new)commmands are
%% removed.
%%
%% Erwin de Gelder, January 6, 2018
%%==========================================================================

%%==========================================================================
%% Title:
%% Author(s):
%% To be published in:
%% Date of submission:
%% Date of submission R1:
%% Date of submission R2:
%% Date of final submission:
%%==========================================================================


%% Generic
\documentclass[10pt,final,a4paper,oneside,onecolumn]{article}

%% IEEE
%% Document class options: replace "draftcls" by "final" for final document.
%% All other options may just as well be omitted because they are the default values.
%\documentclass[10pt,final,journal,letterpaper,twoside,twocolumn]{IEEEtran}


%%==========================================================================
%% Document automation
%%==========================================================================

\def\reptitle{Completeness of real-world driving data}
\def\repauthor{Erwin de Gelder, etc.}


%%==========================================================================
%% Packages
%%==========================================================================

\usepackage[a4paper,left=3.5cm,right=3.5cm,top=3cm,bottom=3cm]{geometry} %% change page layout; remove for IEEE paper format
\usepackage[T1]{fontenc}                        %% output font encoding for international characters (e.g., accented)
\usepackage[cmex10]{amsmath}                    %% math typesetting; consider using the [cmex10] option
\usepackage{amssymb}                            %% special (symbol) fonts for math typesetting
\usepackage{amsthm}                             %% theorem styles
\usepackage{dsfont}                             %% double stroke roman fonts: the real numbers R: $\mathds{R}$
\usepackage{mathrsfs}                           %% formal script fonts: the Laplace transform L: $\mathscr{L}$
\usepackage[pdftex]{graphicx}                   %% graphics control; use dvips for TeXify; use pdftex for PDFTeXify
\usepackage{array}                              %% array functionality (array, tabular)
\usepackage{upgreek}                            %% upright Greek letters; add the prefix 'up', e.g. \upphi

\usepackage[utf8]{inputenc}   				 	%% utf8 support (required for biblatex)
\usepackage[style=ieee,doi=false,isbn=false,url=false,date=year,minbibnames=15,maxbibnames=15,backend=biber]{biblatex}
%\renewcommand*{\bibfont}{\footnotesize}		%% Use this for papers
\setlength{\biblabelsep}{\labelsep}
\bibliography{../bib}

\usepackage{stfloats}                           %% improved handling of floats
\usepackage{multirow}                           %% cells spanning multiple rows in tables
%\usepackage{subfigure}                         %% subfigures and corresponding captions (for use with IEEEconf.cls)
%\usepackage{subfig}                             %% subfigures (IEEEtran.cls: set caption=false)
\usepackage{fancyhdr}                           %% page headers and footers
\usepackage[official,left]{eurosym}             %% the euro symbol; command: \euro
\usepackage{appendix}                           %% appendix layout
\usepackage{xspace}                             %% add space after macro depending on context
\usepackage{verbatim}                           %% provides the comment environment
\usepackage[dutch,USenglish]{babel}             %% language support
\usepackage{wrapfig}                            %% wrapping text around figures
\usepackage{longtable}                          %% tables spanning multiple pages
\usepackage{pgfplots}                           %% support for TikZ figures (Matlab)
\usepackage[breaklinks=true,hidelinks,          %% implement hyperlinks (dvips yields minor problems with breaklinks;
            bookmarksnumbered=true]{hyperref}   %% IEEEtran: set bookmarks=false)
%\usepackage[hyphenbreaks]{breakurl}            %% allow line breaks in URLs (don't use with PDFTeX)
\usepackage{lmodern} 
\usepackage{etoolbox}							%% Needed for apptocmd later
\usepackage[capitalize]{cleveref}
\usepackage{units}
\usepackage{subcaption}
\usepackage{csquotes}							%% Quoted texts are typeset according to rules of main language

%%==========================================================================
%% Fancy headers and footers
%%==========================================================================

\newtoggle{standalone}
\toggletrue{standalone}
\pagestyle{fancy}                                       %% set page style
\fancyhf{}                                              %% clear all header & footer fields
\iftoggle{standalone}{%
	\fancyhead[L]{\includegraphics[width=15mm]{Streetwise}} %% define headers (LE: left field/even pages, etc.)
	\fancyhead[R]{\small\emph{\reptitle}}                   %% similar
	\fancyfoot[C]{\thepage}                                 %% define footer
	\setlength{\headheight}{18pt}                           %% increase head height to accommodate an oversized picture in the header
	\renewcommand*{\headrulewidth}{0.25pt}                  %% header line width
	%% Redefine the default "plain" page style (automatically activated by \maketitle, \section, ...)
	\fancypagestyle{plain}{%
		\fancyhf{}
		\fancyhead[C]{\includegraphics[width=30mm]{Streetwise}}
		\renewcommand{\headrulewidth}{0pt}
		\renewcommand{\footrulewidth}{0pt}
		\setlength{\headheight}{36pt}
	}
}{%
	\renewcommand*{\headrulewidth}{0pt}                  	%% No line in this case
	%% Redefine the default "plain" page style (automatically activated by \maketitle, \section, ...)
	\fancypagestyle{plain}{%
		\fancyhf{}
	}
}
\renewcommand*{\footrulewidth}{0pt}                     %% footer line width


%%==========================================================================
%% TikZ figures
%%==========================================================================

\newlength\figurewidth
\setlength\figurewidth{0.35\textwidth}              %% set figure width
\newlength\figureheight
\setlength\figureheight{0.3\textwidth}              %% set figure height
\pgfplotsset{every axis/.append style={
    scaled y ticks=false,
    scaled x ticks=false,
    y tick label style={/pgf/number format/fixed},
    x tick label style={/pgf/number format/fixed},
    legend style={font=\small}},
    compat=1.9}                                     %% PGFPlots package options
\usetikzlibrary{shapes.geometric, arrows, arrows.meta}
\usepgfplotslibrary{groupplots}
%\usetikzlibrary{external}                           %% Create pdf figures from TikZ. Use PDFTeXify ...
%\tikzexternalize[prefix=./tikz/]                    %% ... with --tex-option=--shell-escape switch.
%\tikzset{external/force remake}                    %% force pdf figure update


%%==========================================================================
%% User-defined commands
%%==========================================================================

\newcommand*{\mat}[1]{\mathbf{#1}}                              %% matrix/vector notation
\newcommand*{\matsym}[1]{\boldsymbol{#1}}                       %% matrix/vector notation for Greek letters
\newcommand*{\T}{^{\scriptscriptstyle\mathsf{T}}}               %% transpose operator
\newcommand*{\Hr}{^{\scriptscriptstyle\mathsf{H}}}              %% conjugate transpose operator
\newcommand*{\ud}{\mathrm{\,d}}                                 %% differential operator (upright d)
\newcommand*{\defeq}{\mathrel{\mathop:}=}                       %% definition sign :=
\newcommand*{\eqdef}{=\mathrel{\mathop:}}                       %% definition sign =:
\newcommand*{\ip}[2]{\left\langle#1\,{,}\,#2\right\rangle}      %% inner product
\newcommand*{\real}[1]{\mathrm{Re}(#1)}                         %% real part
\newcommand*{\imag}[1]{\mathrm{Im}(#1)}                         %% imaginary part
\newcommand*{\lsup}[1]{{}^{#1}\!}                               %% left superscript
\newcommand*{\hi}[1]{$^\text{#1}$}                              %% superscript in normal text
\newcommand*{\lo}[1]{$_\text{#1}$}                              %% subscript in normal text
\newcommand*{\w}[1]{\mathrm{#1}}                                %% multiple character super-/subscript in math mode
\newcommand*{\capskip}{\vspace{-12pt}}                          %% caption skip for figures with subfloats
\newcommand*{\etal}{et al.}                                     %% may be required for Natbib bibliography styles
\renewcommand*{\qedsymbol}{$\blacksquare$}                      %% redefine the end-of-proof symbol
\renewcommand*{\labelitemi}{$\bullet$}                          %% first level item list bullet
\renewcommand*{\labelitemii}{$-$}                               %% second level item list bullet
%\renewcommand*{\theenumi}{\textit{\roman{enumi}}}               %% first level enumerator
\renewcommand*{\labelenumi}{\theenumi.}
\renewcommand*{\theenumii}{\textit{\alph{enumii}}}              %% second level enumerator
\renewcommand*{\labelenumii}{\theenumii.}
\DeclareMathOperator{\tr}{tr}                                   %% trace of a matrix
\DeclareMathOperator{\sgn}{sgn}                                 %% signum function
\DeclareMathOperator{\atan}{atan}                               %% arc tangent

\newcommand{\deuclid}{d_{\textup{Euclidian}}}
\newcommand{\ddtw}{d_{\textup{DTW}}}
\newcommand{\imagi}{\textbf{j}}
\newcommand{\profile}[1]{\textbf{v}_{#1}}
\newcommand{\barprofile}[1]{\bar{\textbf{v}}_{#1}}
\newcommand{\parvstart}{v_{\textup{start}}}
\newcommand{\parvend}{v_{\textup{end}}}
\newcommand{\pardv}{v_{\Delta}}
\newcommand{\pardt}{t_{\textup{duration}}}
\newcommand{\expectation}[1]{\textup{E} \left[ #1 \right]}
\newcommand{\entropy}[1]{\textup{H} \left[ #1 \right]}
\newcommand{\hatentropy}[1]{\hat{\textup{H}} \left[ #1 \right]}
\newcommand{\kl}[2]{\textup{KL} \left( #1 || #2 \right)}

\crefname{figure}{Figure}{Figures}
\crefname{equation}{}{}
\Crefname{equation}{Equation}{Equations}

%%==========================================================================
%% User-defined environments
%%==========================================================================

\theoremstyle{plain}\newtheorem{definition}{Definition}[section]    %% definition
                    \newtheorem{theorem}{Theorem}[section]          %% theorem
                    \newtheorem{lemma}[theorem]{Lemma}              %% lemma
                    \newtheorem{corollary}[theorem]{Corollary}      %% corollary
                    \newtheorem{assumption}{Assumption}[section]    %% assumption
                    \newtheorem{condition}{Condition}[section]      %% condition
\theoremstyle{definition}\newtheorem{example}{Example}[section]     %% examples
\theoremstyle{remark}\newtheorem{remarkenv}{Remark}[section]        %% remarks
\newenvironment{remark}{\begin{remarkenv}}%
                       {\hfill$\blacklozenge$\end{remarkenv}}       %% end remark with a lozenge
\newenvironment{reviewer}{\itshape}{\upshape}                       %% environment for reviewer's comments


%%==========================================================================
%% Miscellaneous
%%==========================================================================

\graphicspath{{./../}{./figures/}{./more_figures/}} %% (graphicx) directory path for figures
%\setlength{\parindent}{0pt}                        %% no paragraph indentation
%\setlength{\parskip}{2ex}                          %% create empty line between paragraphs
%\interdisplaylinepenalty=2500                      %% (amsmath) allow for page breaks within multiline equations
%\numberwithin{equation}{section}                   %% (amsmath) include section number in equation numbering
%\numberwithin{figure}{section}                     %% (amsmath) include section number in figure numbering
%\numberwithin{table}{section}                      %% (amsmath) include section number in table numbering
\addtolength{\arraycolsep}{-0.5mm}                  %% squeeze matrix columns a little
\fnbelowfloat                                       %% (stfloats) put footnote below a float at the page bottom
\urlstyle{same}                                     %% (hyperref) use current font for URLs
\hypersetup{pdftitle={\reptitle},
            pdfauthor={\repauthor}}                 %% (hyperref) pdf properties title and author
%\raggedbottom                                      %% don't add inter-paragraph spacing to achieve \textheight
%\setlength\subfigcapskip{-3pt}                     %% (subfigure) distance between subfloat and subcaption
%\setlength\subfigbottomskip{-3pt}                  %% (subfigure) distance between subcaption and caption
%\renewcommand*{\subcapsize}{\small}                %% (subfigure) subcaption font size
\renewcommand*{\thesubfigure}{(\alph{subfigure})}   %% (subfig) implement 1(a) instead of 1a ...as subfigure reference
\captionsetup[subfloat]{labelformat=simple}        %% ... as subfigure reference
%\captionsetup[subfloat]{
%    farskip=0pt,
%    nearskip=8pt,
%    captionskip=1.5pt,
%    labelfont={small,bf},
%    textfont=small}                                %% (subfig) subfloat caption format
%\captionsetup[figure]{
%    labelfont={small,bf},
%    textfont=small}                                %% (subfig/caption) figure caption format
%\captionsetup[table]{
%    aboveskip=2pt,
%    labelfont={small,bf},
%    textfont={small,sc}}                           %% (subfig/caption) table caption format
\apptocmd{\thebibliography}{\raggedright}{}{}		%% Suppress badness warnings in bibliography


% Table stuff
\usepackage{booktabs}
\usepackage{tabularx}
\setlength{\heavyrulewidth}{0.1em}
\newcommand{\otoprule}{\midrule[\heavyrulewidth]}

\usepackage{changebar}  %% When including the package earlier, it does not work for some reason...

%%==========================================================================
%% Begin document
%%==========================================================================

\begin{document}

\selectlanguage{USenglish}

\title{\textbf{\reptitle}}
\author{\repauthor}
\date{\today}
\maketitle

\tableofcontents

\newpage

\section{Introduction}
\label{sec:introduction}

% Scenario-based approach for assessment of AVs -> this scenarios are important
An important aspect in the development of automated vehicles (AVs) is the assessment of quality and performance aspects of the AVs, such as safety, comfort, and efficiency \cite{bengler2014threedecades, stellet2015taxonomy, wachenfeld2016release, putz2017pegasus, roesener2016scenariobased, kompass2015sicherheitsveranderung}. 
For legal and public acceptance, it is important that there is a clear definition of system performance and that there are quantitative measures for the system quality. 
The more traditional methods \cite{response2006code, ISO26262}, used for evaluation of driver assistance systems, are no longer valid for the assessment of quality and performance aspects of an AV \cite{wachenfeld2016release}. 
Therefore, a scenario-based approach is proposed \cite{roesener2016scenariobased, putz2017pegasus, kompass2015sicherheitsveranderung}. 
For the scenario-based assessment, proper specification of scenarios is crucial since they are directly reflected in test cases used for scenario-based assessment \cite{stellet2015taxonomy}.

% Completeness
To draw conclusions on how an automated vehicle would perform in real-world traffic, it is necessary to know how representative the scenario database, that is used for the scenario-based assessment of the automated vehicle, is. Therefore it is important to quantify how complete the scenario database is \cite{geyer2014, alvarez2017prospective, stellet2015taxonomy}.

\cbstart
To quantify how complete a scenario database is, a few assumptions are made. First, it is assumed that a scenario is a quantitative description, whereas a \emph{scenario class} refers to a qualitative description \cite{degelder2018ontology}. As such, a specific scenario class refers to multiple scenarios. The problem of quantifying the completeness can now be divided into two subproblems.

The first subproblem deals with the quantification of the completeness regarding the scenario classes. To tackle this problem, each scenario class is described by a set of \emph{tags}. For example, the scenario class ``lead vehicle braking at straight road, in daylight, under clear weather conditions'' can be described by the tags ``lead vehicle braking'', ``straight road'', ``daylight'', and ``clear weather''. It is assumed that a scenario class is unambiguously defined by a set of tags. Therefore, the problem of quantifying the completeness regarding the scenario classes boils down to a combinatorial problem.

The second subproblem deals with the quantification of the completeness regarding the scenarios of a specific scenario class. It is assumed that all scenarios of a specific scenario class can be parametrized in a similar manner. Though these parameters can be either discrete or continuous, is it expected that at least one parameter is continuous. As a result, there are infinite number of scenarios that belong to a specific scenario class. Therefore, it is expected that the problem of quantifying the completeness regarding the scenarios of a specific scenario class boils down to a statistical problem. 
\cbend

% Approach
%To achieve a complete database, the process shown in \cref{fig:completeness collection} is adopted. Newly collected data is compared with an existing database. If the data is not ``similar'', the data is added and the process is repeated. If, at some point, the data is ``similar'', it is concluded that the database is complete. This document describes different ways to determine to what degree the newly collected data is ``similar'' to the earlier collected data.

%\tikzstyle{startstop}=[rectangle, rounded corners, minimum width=4em, minimum height=4em, text centered, draw=black, fill=red!30, text width=4em]
%\tikzstyle{process} = [rectangle, minimum width=9em, text width=8.5em, minimum height=4em, text centered, draw=black, fill=orange!30, node distance=10em]
%\tikzstyle{decision} = [diamond, minimum width=4em, text width=4em, minimum height=4em, text centered, draw=black, fill=green!30, node distance=10em]
%\tikzstyle{arrow} = [-{Triangle[width=2mm, length=3mm, round]}, line width=0.6mm] % rounded corners=0.4em
%\begin{figure}[b]
%	\centering
%	\begin{tikzpicture}
%		% Draw the blocks
%		\node[startstop](start){Start};
%		\node[process, right of=start, node distance=7em, text width=5em, minimum width=5.5em](collect){Collect new data};
%		\node[process, right of=collect, node distance=9em](compare){Compare new data with database};
%		\node[decision, right of=compare](similar){New data similar?};
%		\node[startstop, right of=similar, node distance=9em](stop){Database complete};
%		\node[process, below of=compare, node distance=7em](add){Add collected data to database};
%		
%		% Arrows
%		\draw[arrow] (start) -- (collect);
%		\draw[arrow] (collect) -- (compare);
%		\draw[arrow] (compare) -- (similar);
%		\draw[arrow] (similar) -- node[anchor=south]{Yes} (stop);
%		\draw[arrow] (similar) |- node[anchor=north]{No} (add);
%		\draw[arrow] (add) -| (collect);
%	\end{tikzpicture}
%	\caption{Schematic process of collecting data until a certain level of completeness is reached.}
%	\label{fig:completeness collection}
%\end{figure}

% Purpose of this document
This document describes different methods for quantifying the completeness of a dataset. To illustrate the different methods, use is made of artificial data. The advantage of using artificial data is that it is easy to obtain a huge amount of data. The data consists of the speed of a vehicle sampled at $\unit[10]{Hz}$. Furthermore, the acceleration is estimated by filtering the derivative of the speed. A portion of the data is shown in \cref{fig:artificial data}. Note that the data is not very realistic, because in real life, the relative time a vehicle is cruising, i.e., driving at a constant speed, will be higher. This is done for practical reasons, because fewer data needs to be generated to obtain a certain amount of braking and accelerating activities. Therefore, for the purpose of illustrating the different methods, the artificial data suffices.

\setlength\figurewidth{0.8\linewidth}              %% set figure width
\setlength\figureheight{0.3\textwidth}              %% set figure height
\begin{figure}
	\centering
	\begin{minipage}{\linewidth}
		\centering
		% This file was created by matplotlib2tikz v0.6.14.
\begin{tikzpicture}

\definecolor{color0}{rgb}{0.12156862745098,0.466666666666667,0.705882352941177}

\begin{axis}[
xlabel={Time [s]},
ylabel={Speed [m/s]},
xmin=0, xmax=500,
ymin=-1.3007710998959, ymax=27.3161930978138,
width=\figurewidth,
height=\figureheight,
tick align=outside,
tick pos=left,
xmajorgrids,
x grid style={lightgray!92.026143790849673!black},
ymajorgrids,
y grid style={lightgray!92.026143790849673!black}
]
\addplot [line width=2.0pt, color0, forget plot]
table {%
0 0.02495226681389
1 0.41412915559112
2 1.20415249714665
3 2.04559320305136
4 2.96349908784175
5 4.05774866180894
6 5.12866058530594
7 6.33534420532051
8 7.02440786605288
9 7.25058904747158
10 7.35198356916793
11 7.53558457779503
12 7.53558457779503
13 7.53558457779503
14 7.53558457779503
15 7.53558457779503
16 7.53558457779503
17 7.53558457779503
18 7.53558457779503
19 7.53558457779503
20 7.53558457779503
21 7.53558457779503
22 7.53558457779503
23 7.53558457779503
24 7.53558457779503
25 7.53558457779503
26 7.52549984694471
27 7.42982944658301
28 7.37313876723882
29 7.35259530108769
30 7.19363515011014
31 7.12198980592494
32 6.92134867417635
33 6.74889592880281
34 6.54825446130259
35 6.34655984429617
36 6.1292326420585
37 5.94098500435689
38 5.69494203639759
39 5.47554865744759
40 5.26543085890218
41 5.02274943310738
42 4.8156281044786
43 4.57639024257758
44 4.28401370382455
45 4.07446246924165
46 3.83874816162199
47 3.52668706196945
48 3.33329751272105
49 3.08057217745097
50 2.80018403051738
51 2.63291621487606
52 2.44146863581963
53 2.26994966908652
54 2.09016527687301
55 1.91280740001334
56 1.71917366974441
57 1.51941733687503
58 1.32376222611795
59 1.10465280538669
60 0.923368041502252
61 0.762297859725797
62 0.549755191500437
63 0.392930330836383
64 0.274226520723161
65 0.184824745227919
66 0.149237927130451
67 0.0642741344089064
68 0
69 0
70 0
71 0
72 0
73 0
74 0
75 0
76 0
77 0
78 0
79 0
80 0
81 0
82 0
83 0
84 0
85 0
86 0
87 0
88 0
89 0
90 0
91 0
92 0
93 0
94 0
95 0
96 0
97 0
98 0
99 0
100 0
101 0
102 0
103 0.0235510472071443
104 0.588314158964943
105 1.19484288421782
106 1.67410401366337
107 2.06507202894072
108 2.48864779705037
109 2.91059763466345
110 3.34263982019536
111 4.12262318863332
112 5.51238060675081
113 7.07373578415886
114 8.77992636596677
115 10.7249260776893
116 10.4672392154854
117 10.0757932891272
118 9.56215224844886
119 9.14376966264727
120 8.67920735721271
121 8.25774714992156
122 7.84226331300954
123 7.33128031296112
124 6.93576950755434
125 6.68855923520865
126 6.64314153361708
127 6.64314153361708
128 6.64314153361708
129 6.64314153361708
130 6.64314153361708
131 6.64314153361708
132 6.64314153361708
133 6.64314153361708
134 6.64314153361708
135 6.64314153361708
136 6.64314153361708
137 6.64314153361708
138 6.64314153361708
139 6.64314153361708
140 6.77074145955114
141 6.89073386191043
142 7.06366690065455
143 7.18749913095217
144 7.32973394535564
145 7.4975372581904
146 8.01658524200401
147 8.70663554583158
148 9.08290544592541
149 9.28402078519807
150 9.46806007221825
151 9.7223105154507
152 10.1818934142756
153 10.7413384655312
154 11.2984796098625
155 11.9176586938559
156 12.5194171938484
157 13.1673989001897
158 13.8185634610461
159 14.4960619237495
160 15.1351324964136
161 15.9025494632512
162 16.7470706337469
163 17.7725443330002
164 19.0158267752213
165 20.1684419727177
166 21.1562915795093
167 22.0079639517925
168 22.8376271644745
169 24.0096642317662
170 23.4896155489254
171 22.5931385655683
172 21.5635094835419
173 20.5878943722396
174 19.6332161485003
175 19.3518373981516
176 19.3518373981516
177 19.3518373981516
178 19.3518373981516
179 19.3518373981516
180 19.3518373981516
181 19.3518373981516
182 19.3518373981516
183 19.3518373981516
184 19.3518373981516
185 19.3518373981516
186 19.3518373981516
187 19.3518373981516
188 19.4025552783945
189 20.2546208335804
190 21.7480743117486
191 23.5130673481488
192 24.896501678608
193 25.5605191776938
194 26.0154219979179
195 26.0154219979179
196 26.0154219979179
197 26.0154219979179
198 26.0154219979179
199 26.0154219979179
200 26.0154219979179
201 26.0154219979179
202 26.0154219979179
203 26.0154219979179
204 26.0154219979179
205 26.0154219979179
206 26.0154219979179
207 26.0154219979179
208 25.7111559691589
209 25.0607980681135
210 24.4066165406515
211 23.6002186510014
212 22.7977079879777
213 22.0465377827481
214 21.5102868770601
215 20.8941465889376
216 20.4073136737359
217 19.9737336244629
218 19.7531383622051
219 19.6295287448304
220 19.6295287448304
221 19.6295287448304
222 19.6295287448304
223 19.6295287448304
224 19.6295287448304
225 19.6295287448304
226 19.6295287448304
227 19.6295287448304
228 19.6295287448304
229 19.6295287448304
230 19.6295287448304
231 19.5715018533823
232 19.0468180577465
233 18.5350770416726
234 18.0998583315973
235 17.7876979316757
236 17.3914555497333
237 17.1015741384041
238 16.7190056305012
239 16.3919276175824
240 16.1888654412743
241 16.1888654412743
242 16.1888654412743
243 16.1888654412743
244 16.1888654412743
245 16.1888654412743
246 16.1888654412743
247 16.1888654412743
248 16.1888654412743
249 16.1888654412743
250 16.1888654412743
251 16.1888654412743
252 16.1888654412743
253 16.1888654412743
254 16.1888654412743
255 16.1888654412743
256 16.1888654412743
257 16.244759975667
258 16.8376672865179
259 17.5051018539717
260 18.2576967555946
261 18.9737845229909
262 19.0598339898966
263 19.3737393645219
264 19.7780051732413
265 19.7750646496584
266 19.5818271855058
267 19.528532939132
268 19.3897714092076
269 19.1730484775524
270 19.0513159615543
271 18.8363440657725
272 18.6405685667133
273 18.4525606006436
274 18.1930288389244
275 18.0740734911633
276 17.7522965652754
277 17.5547049457927
278 17.3370619938893
279 16.9285281165711
280 16.7876828455576
281 16.4786935853557
282 16.2479875656713
283 15.9129409464543
284 15.5947025134267
285 15.3516940963531
286 14.9490347618845
287 14.7103429485559
288 14.3484474459726
289 14.0220341010075
290 13.7335911333508
291 13.3242325132686
292 13.0725891327197
293 12.7156900135848
294 12.3962597387869
295 12.0950710845955
296 11.8470410830888
297 11.4702597792302
298 11.1122290580795
299 10.8289092166501
300 10.4307816182326
301 10.231867081881
302 9.79736215562572
303 9.48660125092045
304 9.18812089915069
305 8.80291457171979
306 8.53319665952585
307 8.18008216362833
308 7.88378830262853
309 7.53092592243636
310 7.18615042641893
311 6.91527427402575
312 6.53860577725202
313 6.28159780446968
314 5.91660645755556
315 5.5632972742191
316 5.30163674433125
317 4.97631348956604
318 4.73939539813799
319 4.43207675746579
320 4.13442934579078
321 3.92553464856429
322 3.63663116278089
323 3.47398684094807
324 3.20007407439704
325 2.96601681879235
326 2.7549204602362
327 2.48956742290133
328 2.31407933796254
329 2.05084871091424
330 1.8603454909734
331 1.63068983844745
332 1.47089449479967
333 1.1962613977489
334 1.00149028199273
335 0.778947418938918
336 0
337 0
338 0
339 0
340 0
341 0
342 0
343 0
344 0
345 0
346 0
347 0
348 0
349 0
350 0
351 0
352 0
353 0
354 0
355 0
356 0
357 0
358 0
359 0
360 0
361 0
362 0
363 0
364 0
365 0
366 0
367 0
368 0
369 0
370 0
371 0
372 0
373 0
374 0
375 0
376 0
377 0
378 0
379 0
380 0
381 0
382 0
383 0
384 0
385 0
386 0
387 0.00892171360271377
388 0.0981386339096986
389 0.362222895776485
390 0.806968384180307
391 1.33290797599514
392 1.71118099526973
393 1.98640461852989
394 2.21926758459486
395 2.37940502248795
396 2.79428684471596
397 2.8749967233432
398 2.89999941338979
399 3.03917831746172
400 3.34429269924698
401 3.61685873113329
402 3.82337750933905
403 4.09145640525861
404 4.3662212620465
405 4.87253283260646
406 5.48946755359851
407 6.01672383070859
408 6.55246704472216
409 7.18725031364922
410 7.81402483661992
411 8.4644207015304
412 9.03584991060398
413 9.57873477076918
414 9.99359583305322
415 10.4066656896986
416 10.7448022582683
417 11.0307436864592
418 10.903825648276
419 10.1346057409828
420 9.1460041169704
421 8.16759985745562
422 7.51621256146151
423 7.51621256146151
424 7.51621256146151
425 7.51621256146151
426 7.51621256146151
427 7.51621256146151
428 7.51621256146151
429 7.51621256146151
430 7.51621256146151
431 7.51621256146151
432 7.51621256146151
433 7.51621256146151
434 7.51621256146151
435 7.41353993523616
436 6.66263204457012
437 5.75267170116231
438 4.86183333207651
439 3.91664049713349
440 3.00768896969473
441 2.03631660135979
442 1.02672437477843
443 0
444 0
445 0
446 0
447 0
448 0
449 0
450 0
451 0
452 0
453 0
454 0
455 0
456 0
457 0
458 0.478592642672584
459 0.770708092597168
460 1.03872527629063
461 1.63218507114834
462 2.36813306141615
463 3.09204464323763
464 3.52951065764046
465 3.55574416565908
466 3.55574416565908
467 3.55574416565908
468 3.55574416565908
469 3.53129014676358
470 3.19387417395677
471 2.77685196794471
472 2.44515921876903
473 2.01200601908199
474 1.61865749700049
475 1.23571423657206
476 0.858497821150841
477 0.50495344521193
478 0.208380993476534
479 0.036681390349861
480 0
481 0
482 0
483 0
484 0
485 0
486 0
487 0
488 0
489 0
490 0
491 0
492 0
493 0
494 0
495 0
496 0
497 0
498 0
499 0
};
\addplot [semithick, black, forget plot]
table {%
10.9 -1.3007710998959
10.9 27.3161930978138
};
\addplot [semithick, black, forget plot]
table {%
25.8 -1.3007710998959
25.8 27.3161930978138
};
\addplot [semithick, black, forget plot]
table {%
67.9 -1.3007710998959
67.9 27.3161930978138
};
\addplot [semithick, black, forget plot]
table {%
102.8 -1.3007710998959
102.8 27.3161930978138
};
\addplot [semithick, black, forget plot]
table {%
115 -1.3007710998959
115 27.3161930978138
};
\addplot [semithick, black, forget plot]
table {%
125.2 -1.3007710998959
125.2 27.3161930978138
};
\addplot [semithick, black, forget plot]
table {%
139.1 -1.3007710998959
139.1 27.3161930978138
};
\addplot [semithick, black, forget plot]
table {%
169 -1.3007710998959
169 27.3161930978138
};
\addplot [semithick, black, forget plot]
table {%
174.5 -1.3007710998959
174.5 27.3161930978138
};
\addplot [semithick, black, forget plot]
table {%
187.3 -1.3007710998959
187.3 27.3161930978138
};
\addplot [semithick, black, forget plot]
table {%
193.5 -1.3007710998959
193.5 27.3161930978138
};
\addplot [semithick, black, forget plot]
table {%
207 -1.3007710998959
207 27.3161930978138
};
\addplot [semithick, black, forget plot]
table {%
218.6 -1.3007710998959
218.6 27.3161930978138
};
\addplot [semithick, black, forget plot]
table {%
230.8 -1.3007710998959
230.8 27.3161930978138
};
\addplot [semithick, black, forget plot]
table {%
239.7 -1.3007710998959
239.7 27.3161930978138
};
\addplot [semithick, black, forget plot]
table {%
256.3 -1.3007710998959
256.3 27.3161930978138
};
\addplot [semithick, black, forget plot]
table {%
264.2 -1.3007710998959
264.2 27.3161930978138
};
\addplot [semithick, black, forget plot]
table {%
336 -1.3007710998959
336 27.3161930978138
};
\addplot [semithick, black, forget plot]
table {%
386.8 -1.3007710998959
386.8 27.3161930978138
};
\addplot [semithick, black, forget plot]
table {%
417.4 -1.3007710998959
417.4 27.3161930978138
};
\addplot [semithick, black, forget plot]
table {%
421.8 -1.3007710998959
421.8 27.3161930978138
};
\addplot [semithick, black, forget plot]
table {%
434.7 -1.3007710998959
434.7 27.3161930978138
};
\addplot [semithick, black, forget plot]
table {%
442.4 -1.3007710998959
442.4 27.3161930978138
};
\addplot [semithick, black, forget plot]
table {%
457.3 -1.3007710998959
457.3 27.3161930978138
};
\addplot [semithick, black, forget plot]
table {%
464.2 -1.3007710998959
464.2 27.3161930978138
};
\addplot [semithick, black, forget plot]
table {%
468.8 -1.3007710998959
468.8 27.3161930978138
};
\addplot [semithick, black, forget plot]
table {%
479.1 -1.3007710998959
479.1 27.3161930978138
};
\end{axis}

\end{tikzpicture}%
		\subcaption{Speed.}	
	\end{minipage}
	\begin{minipage}{\linewidth}
		\centering
		% This file was created by matplotlib2tikz v0.6.14.
\begin{tikzpicture}

\definecolor{color0}{rgb}{0.12156862745098,0.466666666666667,0.705882352941177}

\begin{axis}[
xlabel={Time [s]},
ylabel={Acceleration [m/s$^2$]},
xmin=0, xmax=500,
ymin=-0.583867892100671, ymax=1.03365861056101,
width=\figurewidth,
height=\figureheight,
tick align=outside,
tick pos=left,
xmajorgrids,
x grid style={white!69.019607843137251!black},
ymajorgrids,
y grid style={white!69.019607843137251!black}
]
\addplot [line width=2.0pt, color0, forget plot]
table {%
0 0.046071478645592
1 0.223788324152102
2 0.421331638033323
3 0.464798538112757
4 0.482002249409247
5 0.542704025554155
6 0.574399055121095
7 0.542666924446114
8 0.13420489500123
9 0.0953488444427908
10 0.0761750555429388
11 0.0550558691324774
12 0
13 0
14 0
15 0
16 0
17 0
18 0
19 0
20 0
21 0
22 0
23 0
24 0
25 0
26 -0.0352965579761229
27 -0.0904236207108999
28 0.0445722714711838
29 -0.0820783750607061
30 -0.0367913433023901
31 -0.0601639867483343
32 -0.0784493977767013
33 -0.134462456073953
34 -0.0842827165247129
35 -0.0927262473973367
36 -0.126747513852968
37 -0.0818528105697847
38 -0.123647925058989
39 -0.13446307800428
40 -0.0855181663174696
41 -0.123081904876108
42 -0.116704283094227
43 -0.108441644972883
44 -0.125511810830222
45 -0.108988397554653
46 -0.111844135434368
47 -0.162061967986071
48 -0.109413235437708
49 -0.110649932503195
50 -0.109472222154236
51 -0.114333773388576
52 -0.102184243617721
53 -0.0275402667619038
54 -0.0703900642035362
55 -0.115351690285978
56 -0.0703906517687086
57 -0.13085770465497
58 -0.0847698096949034
59 -0.0886543480692636
60 -0.114052644316111
61 -0.0942297514652619
62 -0.0492804537803637
63 -0.100847129921372
64 -0.0470028489283713
65 -0.0179682445708618
66 -0.0650006786147294
67 -0.0168071383342656
68 -0.0151269469870342
69 0
70 0
71 0
72 0
73 0
74 0
75 0
76 0
77 0
78 0
79 0
80 0
81 0
82 0
83 0
84 0
85 0
86 0
87 0
88 0
89 0
90 0
91 0
92 0
93 0
94 0
95 0
96 0
97 0
98 0
99 0
100 0
101 0
102 0
103 0.134055989309501
104 0.291547260042344
105 0.249141569148859
106 0.22211138286595
107 0.191194368103815
108 0.20012622789391
109 0.204477647463156
110 0.259529361446484
111 0.589893743081626
112 0.752653073680113
113 0.824874267458022
114 0.803756772596651
115 0.775016650592568
116 -0.163967803419968
117 -0.216213049124093
118 -0.21621304912365
119 -0.226028777388796
120 -0.226797993283258
121 -0.188215550766635
122 -0.221039569816175
123 -0.252702669714406
124 -0.180626613636798
125 -0.0899692925870035
126 0
127 0
128 0
129 0
130 0
131 0
132 0
133 0
134 0
135 0
136 0
137 0
138 0
139 0.0112666493814668
140 0.0976211146543804
141 0.0821237532968198
142 0.0574162973926828
143 0.0533929425846384
144 0.0945309760158537
145 0.183042090733167
146 0.315517819870861
147 0.239519357537024
148 0.104498064216776
149 0.085440460337173
150 0.0922274314092482
151 0.16728303653167
152 0.313519610251634
153 0.29959475470192
154 0.258423489162658
155 0.276204328951993
156 0.31767247068215
157 0.322393310489969
158 0.326042289921642
159 0.337235629420189
160 0.369909406826733
161 0.421391243998626
162 0.435082977463464
163 0.58539897645362
164 0.616917177656438
165 0.545976504096689
166 0.430030468509191
167 0.421351302736131
168 0.386211437278508
169 0.542982746437199
170 -0.423115386130016
171 -0.505602817397868
172 -0.497019231751894
173 -0.488853954669262
174 -0.38919913516235
175 0
176 0
177 0
178 0
179 0
180 0
181 0
182 0
183 0
184 0
185 0
186 0
187 0
188 0.172181783824227
189 0.527468553480686
190 0.960134678621845
191 0.874693575766203
192 0.417449559780527
193 0.243057302678519
194 0
195 0
196 0
197 0
198 0
199 0
200 0
201 0
202 0
203 0
204 0
205 0
206 0
207 -0.0199676556226294
208 -0.28429933967092
209 -0.485410778709468
210 -0.330906695270528
211 -0.320495884621208
212 -0.273822298396261
213 -0.305194967311134
214 -0.232948452849906
215 -0.258607535529459
216 -0.269549605525112
217 -0.178750273717176
218 -0.0684604405391145
219 0
220 0
221 0
222 0
223 0
224 0
225 0
226 0
227 0
228 0
229 0
230 0
231 -0.196023484621696
232 -0.184981044197633
233 -0.257286564272734
234 -0.209476629906934
235 -0.177783117183525
236 -0.127538911484843
237 -0.205996709777201
238 -0.246991402302909
239 -0.171840515835353
240 -0.018498189804065
241 0
242 0
243 0
244 0
245 0
246 0
247 0
248 0
249 0
250 0
251 0
252 0
253 0
254 0
255 0
256 0
257 0.198388835082422
258 0.318909331479071
259 0.366468133526771
260 0.359237910983561
261 0.17947442864406
262 0.0213312552743101
263 0.283088592489635
264 0.030136975180179
265 -0.0382790035960223
266 -0.053841528432617
267 -0.0468701674654302
268 -0.0963017948611586
269 -0.107287959573718
270 -0.0766475499485946
271 -0.114248089273097
272 -0.0674112564024316
273 -0.116284560976727
274 -0.101552622778108
275 -0.128148315943804
276 -0.154148468286547
277 -0.0385207179442464
278 -0.159686121378767
279 -0.128459823641437
280 -0.110883355630998
281 -0.211861857800692
282 -0.0596666964618748
283 -0.190100796309232
284 -0.0896155882043139
285 -0.231124307665471
286 -0.173275577070223
287 -0.122514683856677
288 -0.186413453465031
289 -0.10927893140022
290 -0.192601688908625
291 -0.19867739337622
292 -0.139295895692826
293 -0.198344575201245
294 -0.0866707600093752
295 -0.19816804808004
296 -0.129737776940262
297 -0.213337423514281
298 -0.0778979437441887
299 -0.211863948694551
300 -0.169549108860421
301 -0.159371900146791
302 -0.203768814750679
303 -0.0866716153745504
304 -0.192601833147396
305 -0.164140246923865
306 -0.163359977715133
307 -0.183036634700208
308 -0.115762456431621
309 -0.1942854496333
310 -0.136043779228423
311 -0.189365750750484
312 -0.182973410235165
313 -0.109791115330625
314 -0.197439774267549
315 -0.117348168401494
316 -0.173343230749103
317 -0.152165405728093
318 -0.122788368028338
319 -0.174898667548872
320 -0.0679314621942386
321 -0.144452692290919
322 -0.128859803700888
323 -0.106015473379167
324 -0.143697508915754
325 -0.0770406755634505
326 -0.125936324425861
327 -0.110265951564871
328 -0.123001879721851
329 -0.125192333318797
330 -0.0676408781789105
331 -0.11613118451899
332 -0.11107408496279
333 -0.144058605564379
334 -0.0385207179442447
335 -0.191505709964151
336 -0.221494128179388
337 0
338 0
339 0
340 0
341 0
342 0
343 0
344 0
345 0
346 0
347 0
348 0
349 0
350 0
351 0
352 0
353 0
354 0
355 0
356 0
357 0
358 0
359 0
360 0
361 0
362 0
363 0
364 0
365 0
366 0
367 0
368 0
369 0
370 0
371 0
372 0
373 0
374 0
375 0
376 0
377 0
378 0
379 0
380 0
381 0
382 0
383 0
384 0
385 0
386 0
387 0.0312259976095071
388 0.114199898120503
389 0.217244190434968
390 0.235976193689842
391 0.31359082060422
392 0.152557446269393
393 -0.0223096012392903
394 0.177529111802229
395 0.0762920967801017
396 0.16436422644095
397 -0.0133942385710046
398 0.114636860328785
399 0.179111987185835
400 0.13160037675178
401 0.201838188945132
402 0.0793710729671111
403 0.136505284928309
404 0.249350148291786
405 0.268083776375465
406 0.231533181485182
407 0.319632657995126
408 0.360879325394285
409 0.353762909380075
410 0.258295434501743
411 0.309593419737499
412 0.294192697721561
413 0.245349801559558
414 0.261403579947372
415 0.298651761408779
416 0.253821565018001
417 0.0477309658798859
418 -0.243936567038888
419 -0.476730095104475
420 -0.486716762820994
421 -0.497981600312129
422 -0.0121824432895057
423 0
424 0
425 0
426 0
427 0
428 0
429 0
430 0
431 0
432 0
433 0
434 0
435 -0.220444947028303
436 -0.476103279737743
437 -0.490968549967181
438 -0.429061393169296
439 -0.461519852038074
440 -0.473859108833874
441 -0.510343960161503
442 -0.472848390469481
443 0
444 0
445 0
446 0
447 0
448 0
449 0
450 0
451 0
452 0
453 0
454 0
455 0
456 0
457 0
458 0.304343765381447
459 0.146409572154353
460 0.193212626652152
461 0.351497430158638
462 0.351497467053642
463 0.324731768491022
464 0.119824626282649
465 0
466 0
467 0
468 0
469 -0.0861100004038642
470 -0.17813859316677
471 -0.171959784957742
472 -0.219957557972776
473 -0.222039165910189
474 -0.196673633308648
475 -0.208899450183817
476 -0.149065012050385
477 -0.145164235491094
478 -0.144513930467346
479 -0.0545016339276104
480 0
481 0
482 0
483 0
484 0
485 0
486 0
487 0
488 0
489 0
490 0
491 0
492 0
493 0
494 0
495 0
496 0
497 0
498 0
499 0
};
\addplot [semithick, black, forget plot]
table {%
10.9 -0.583867892100671
10.9 1.03365861056101
};
\addplot [semithick, black, forget plot]
table {%
25.8 -0.583867892100671
25.8 1.03365861056101
};
\addplot [semithick, black, forget plot]
table {%
67.9 -0.583867892100671
67.9 1.03365861056101
};
\addplot [semithick, black, forget plot]
table {%
102.8 -0.583867892100671
102.8 1.03365861056101
};
\addplot [semithick, black, forget plot]
table {%
115 -0.583867892100671
115 1.03365861056101
};
\addplot [semithick, black, forget plot]
table {%
125.2 -0.583867892100671
125.2 1.03365861056101
};
\addplot [semithick, black, forget plot]
table {%
139.1 -0.583867892100671
139.1 1.03365861056101
};
\addplot [semithick, black, forget plot]
table {%
169 -0.583867892100671
169 1.03365861056101
};
\addplot [semithick, black, forget plot]
table {%
174.5 -0.583867892100671
174.5 1.03365861056101
};
\addplot [semithick, black, forget plot]
table {%
187.3 -0.583867892100671
187.3 1.03365861056101
};
\addplot [semithick, black, forget plot]
table {%
193.5 -0.583867892100671
193.5 1.03365861056101
};
\addplot [semithick, black, forget plot]
table {%
207 -0.583867892100671
207 1.03365861056101
};
\addplot [semithick, black, forget plot]
table {%
218.6 -0.583867892100671
218.6 1.03365861056101
};
\addplot [semithick, black, forget plot]
table {%
230.8 -0.583867892100671
230.8 1.03365861056101
};
\addplot [semithick, black, forget plot]
table {%
239.7 -0.583867892100671
239.7 1.03365861056101
};
\addplot [semithick, black, forget plot]
table {%
256.3 -0.583867892100671
256.3 1.03365861056101
};
\addplot [semithick, black, forget plot]
table {%
264.2 -0.583867892100671
264.2 1.03365861056101
};
\addplot [semithick, black, forget plot]
table {%
336 -0.583867892100671
336 1.03365861056101
};
\addplot [semithick, black, forget plot]
table {%
386.8 -0.583867892100671
386.8 1.03365861056101
};
\addplot [semithick, black, forget plot]
table {%
417.4 -0.583867892100671
417.4 1.03365861056101
};
\addplot [semithick, black, forget plot]
table {%
421.8 -0.583867892100671
421.8 1.03365861056101
};
\addplot [semithick, black, forget plot]
table {%
434.7 -0.583867892100671
434.7 1.03365861056101
};
\addplot [semithick, black, forget plot]
table {%
442.4 -0.583867892100671
442.4 1.03365861056101
};
\addplot [semithick, black, forget plot]
table {%
457.3 -0.583867892100671
457.3 1.03365861056101
};
\addplot [semithick, black, forget plot]
table {%
464.2 -0.583867892100671
464.2 1.03365861056101
};
\addplot [semithick, black, forget plot]
table {%
468.8 -0.583867892100671
468.8 1.03365861056101
};
\addplot [semithick, black, forget plot]
table {%
479.1 -0.583867892100671
479.1 1.03365861056101
};
\end{axis}

\end{tikzpicture}%
		\subcaption{Acceleration.}	
	\end{minipage}
	\caption{Artificial data of the (a) speed and (b) acceleration of a vehicle. The black vertical lines indicate the events. In between two events, the vehicle is either acceleration, cruising, or braking.}
	\label{fig:artificial data}	
\end{figure}

% Structure of document
The methods for estimating whether the amount of driving data is sufficient can be divided into two groups, i.e., the sample-based methods and the activity-based methods. The sample-based methods are discussed in \cref{sec:sample based} and the activity-based methods are discussed in \cref{sec:activity based methods}.


\section{Sample-based methods}
\label{sec:sample based}


\subsection{Comparing probability density functions}
\label{sec:sample based pdf}

According to Wang \etal\ \cite{wang2017much}, the question of ``how much driving data is sufficient to address problems such as the cause of accidents, distraction and inattention, eco-driving styles, modeling driver behavior, and the effects of driver assistance systems on driver behavior'' has not been fully studied. Therefore, Wang \etal\ \cite{wang2017much} present a method for quantifying the completeness regarding naturalistic driving data. 

Wang \etal\ \cite{wang2017much} use a statistical method for determining the appropriate amount of data. Based on the data, a probability density function (PDF) is estimated. The point at which the PDF converges determines the amount of appropriate data. 

Let the PDF of a certain state $x \in \mathds{R}^D$ be denoted by $f(x)$. Based on the $n$ data points $\textbf{x}=\{x_i\}_{i=1}^n$, the PDF $f(x)$ can be estimated. This estimation is denoted by $\hat{f}(x;n)$. For estimating $f(x)$, Kernel Density Estimation (KDE) \cite{rosenblatt1956remarks, parzen1962estimation} is employed, i.e.,
\begin{equation} \label{eq:kde}
	\hat{f}(x;n) = \sum_{i=1}^n \frac{1}{h^D n} K \left( \frac{x - x_i}{h} \right),
\end{equation}
where $h$ denotes the so-called bandwidth and $K(\cdot)$ represents the kernel. A Gaussian kernel is used, i.e.,
\begin{equation}
	K(x) = \frac{1}{\sqrt{2\pi}} \exp \left\{ -\frac{||x||^2}{2}\right\}.
\end{equation}
For the bandwidth, the Silverman rule \cite{silverman1986density} is used, i.e., 
\begin{equation} \label{eq:silverman}
	h=1.06\cdot\hat{\sigma}\cdot n^{-1/5},
\end{equation}
where $\hat{\sigma}$ is the standard deviation of the data $\{x_i\}_{i=1}^n$.

\begin{remark}
	The power of the Silverman rule is in its simplicity. The disadvantage, however, is that it tends to oversmooth the density when it is used for multimodal and highly skewed densities \cite{silverman1986density}. When comparing two estimated PDFs, they will be more similar when both are oversmoothed. As the comparison of two estimated PDFs leads to the quantification of completeness, the use of the Silverman rule might lead to an overestimation of the degree of completeness.
\end{remark}

The distribution $f(\cdot;n)$ will change as $n$ increases. It is expected, however, that the change will decrease as $n$ increases. Therefore, the distribution $\hat{f}(\cdot;n+m)$ with $m \in \mathds{R}^+$ is compared with $\hat{f}(\cdot;n)$. For the comparison of two PDFs, the Kullback-Leibler divergence is employed \cite{kullback1951, bishop2006pattern}. For the two distributions $\hat{f}(\cdot;n+m)$ and $\hat{f}(\cdot;n)$, the Kullback-Leibler divergence is as follows:
\begin{equation} \label{eq:KL divergence}
	\kl{\hat{f}(x;n+m)}{\hat{f}(x;n)} = \int_{-\infty}^{\infty} \hat{f}(x;n+m) \ln \left( \frac{\hat{f}(x;n+m)}{\hat{f}(x;n)} \right) \ud x.
\end{equation}

The Kullback-Leibler divergence is always positive and zero if and only if $\hat{f}(\cdot;n+m) = \hat{f}(\cdot;n)$. Wang~\etal\ \cite{wang2017much} determine the proper amount of driving data such that $\kl{\hat{f}(\cdot;n+m)}{\hat{f}(\cdot;n)}$ changes very slightly if more data samples are to be added, i.e.,
\begin{equation} \label{eq:completeness wang}
	\left| \kl{\hat{f}(\cdot;n+m)}{\hat{f}(\cdot;n)} - \kl{\hat{f}(\cdot;n+2m)}{\hat{f}(\cdot;n+m)} \right| \leq \epsilon,
\end{equation}
for a small value of $\epsilon\in\mathds{R}^+$. The proper number of samples $n$ equals the minimum value of $n$ for which \cref{eq:completeness wang} is satisfied.

\begin{remark}
	Although Wang~\etal\ use \cref{eq:completeness wang} to determine the appropriate number of samples $n$, it seems a bad choice as \cref{eq:completeness wang} is satisfied if the difference between $\hat{f}(\cdot;n+m)$ and $\hat{f}(\cdot;n)$ is almost similar to the difference between $\hat{f}(\cdot;n+m)$ and $\hat{f}(\cdot;n+2m)$. It would be more appropriate to use
	\begin{equation} \label{eq:kl epsilon}
		\kl{\hat{f}(\cdot;n+m)}{\hat{f}(\cdot;n)} \leq \epsilon,
	\end{equation}
	for a small value of $\epsilon\in\mathds{R}^+$.
	
	From the work of Wang~\etal, \cref{eq:kl epsilon} is not considered. Furthermore, no argumentation is provided for choosing \cref{eq:completeness wang} as condition for the completeness of the driving data.
\end{remark}

\begin{example}
	\Cref{eq:completeness wang} is applied for the artificial data as shown in \cref{fig:artificial data}. When using similar values for $m$ and $\epsilon$ as in \cite{wang2017much}, i.e., $m=2000$ and $\epsilon=10^{-4}$, and the velocity as the feature $x$, it appears that 27 minutes of artificial data suffices. The Kullback-Leibler divergence \cref{eq:KL divergence} is shown in \cref{fig:kl sample based} as a function of the amount of data. \Cref{fig:kl sample based} shows that as $n$ increases, the Kullback-Leibler divergence becomes less. 
	
	\setlength\figurewidth{0.6\linewidth}
	\setlength\figureheight{0.7\figurewidth}
	\begin{figure}
		\centering
		% This file was created by matplotlib2tikz v0.6.14.
\begin{tikzpicture}

\begin{axis}[
xlabel={Time [min]},
ylabel={KL divergence $\times$ $10^3$},
xmin=-10.5, xmax=293.833333333333,
ymin=0, ymax=10,
width=\figurewidth,
height=\figureheight,
tick align=outside,
tick pos=left,
xmajorgrids,
x grid style={lightgray!92.026143790849673!black},
ymajorgrids,
y grid style={lightgray!92.026143790849673!black},
legend entries={{Speed},{Acceleration}},
legend cell align={left},
legend style={draw=white!80.0!black}
]
\addlegendimage{no markers, blue}
\addlegendimage{no markers, red}
\addplot [semithick, blue]
table {%
3.33333333333333 85.5559705975772
6.66666666666667 10.597241997071
10 115.017598250175
13.3333333333333 1064.05436314479
16.6666666666667 10.6017651787707
20 4.45501288529982
23.3333333333333 6.49593773314267
26.6666666666667 5.38107335114233
30 6.22951068560297
33.3333333333333 2.08797269193545
36.6666666666667 3.88246066550233
40 29.9544288763864
43.3333333333333 12.8342313126479
46.6666666666667 1.59204298109914
50 3.4036252603755
53.3333333333333 1.08962972854284
56.6666666666667 2.50556811561683
60 1.3585016097844
63.3333333333333 0.531319059821967
66.6666666666667 1.05755232810577
70 0.418669519318776
73.3333333333333 2.13083579658864
76.6666666666667 0.987240093555623
80 2.75234155777671
83.3333333333333 1.82257544000031
86.6666666666667 1.77411962237324
90 0.212710491081484
93.3333333333333 0.943797802243864
96.6666666666667 0.299340423297768
100 2.00766393867795
103.333333333333 0.595403608095591
106.666666666667 0.831800148919101
110 0.495944553052386
113.333333333333 0.270618319363876
116.666666666667 1.12177306642093
120 0.282825727089753
123.333333333333 0.506209465093898
126.666666666667 0.17453847528837
130 1.31645943069852
133.333333333333 0.173995438563436
136.666666666667 0.322104475259102
140 0.728212850822317
143.333333333333 0.128255224219095
146.666666666667 0.119871938643659
150 0.983017170580136
153.333333333333 0.909759198317718
156.666666666667 1.93890221026825
160 0.860775689799353
163.333333333333 0.163793851055491
166.666666666667 0.253018193456412
170 0.286440926288585
173.333333333333 0.22242266248924
176.666666666667 0.176239731428782
180 0.164983494887799
183.333333333333 0.200264542658002
186.666666666667 0.313149602005561
190 0.0906603054456458
193.333333333333 0.431474795628275
196.666666666667 0.293819722115019
200 0.280105660116808
203.333333333333 0.16291867260442
206.666666666667 0.163020621390103
210 0.127210963511037
213.333333333333 0.147487535508811
216.666666666667 0.0953089234725422
220 0.411792780306419
223.333333333333 0.160783488582505
226.666666666667 0.0879132224508249
230 0.0285110422529281
233.333333333333 0.0518486994991802
236.666666666667 0.102624107219071
240 0.153335176031456
243.333333333333 0.0582418167355966
246.666666666667 0.41487242287834
250 0.197415703503801
253.333333333333 0.182695430229512
256.666666666667 0.114067715296794
260 0.0878856303241455
263.333333333333 0.0384522887079879
266.666666666667 0.0262138920313726
270 0.0763982673903383
273.333333333333 0.0566523556784322
276.666666666667 0.172680670425977
280 0.146711714787699
};
\addplot [semithick, red, dashed]
table {%
3.33333333333333 59.2993740411287
6.66666666666667 77.1754349322197
10 33.1066469247452
13.3333333333333 inf
16.6666666666667 4.58573209363791
20 9.01710684721958
23.3333333333333 1.89199289576062
26.6666666666667 2.81288840310293
30 2.34227675855516
33.3333333333333 3.76309225233826
36.6666666666667 3.74515580645727
40 1.44975668683986
43.3333333333333 0.707569653516197
46.6666666666667 1.40978765093012
50 1.00608068726325
53.3333333333333 0.615700265466338
56.6666666666667 1.57989076410248
60 0.770026812072535
63.3333333333333 0.243033319522599
66.6666666666667 0.545102872306463
70 2.52850695134444
73.3333333333333 0.225439433018832
76.6666666666667 0.470045393326763
80 0.82805708796364
83.3333333333333 0.873852997126415
86.6666666666667 0.250293025332482
90 0.509160704916874
93.3333333333333 0.577001783286116
96.6666666666667 0.279467153020571
100 0.139130783967779
103.333333333333 0.327010105133279
106.666666666667 0.508294759863066
110 0.326735429741867
113.333333333333 0.29494049906929
116.666666666667 0.279275946486387
120 1.22481206594918
123.333333333333 0.336863149067572
126.666666666667 0.328662173576838
130 0.30500939238111
133.333333333333 0.337186109757032
136.666666666667 0.221493885405066
140 0.171818450677656
143.333333333333 0.160308351683289
146.666666666667 0.147859811079427
150 0.266061848705773
153.333333333333 0.119649656069875
156.666666666667 0.320645489680288
160 0.249546146072053
163.333333333333 0.247416239969758
166.666666666667 0.359469608346775
170 0.440149950539897
173.333333333333 0.401371132493005
176.666666666667 0.397189691608965
180 0.233600339981118
183.333333333333 0.126581525501493
186.666666666667 0.157270453690798
190 0.0969993616583276
193.333333333333 0.238320496575932
196.666666666667 0.243881151976801
200 0.115638498385969
203.333333333333 0.330808644375143
206.666666666667 0.113000252909512
210 0.0922638353675995
213.333333333333 0.358646519403488
216.666666666667 0.671665577776142
220 0.315959861398592
223.333333333333 -0.000326424980687322
226.666666666667 0.352587058419569
230 -0.0141259465440095
233.333333333333 0.263515301271314
236.666666666667 0.205854346062891
240 0.338505644398863
243.333333333333 0.884167954698448
246.666666666667 0.153825887710452
250 0.0613495057940406
253.333333333333 0.149011363809117
256.666666666667 -0.19050917980561
260 0.154568906158475
263.333333333333 0.00532269676724224
266.666666666667 inf
270 0.0987230314005247
273.333333333333 -0.0790065167667762
276.666666666667 -0.0379722599322027
280 0.282535571753087
};
\end{axis}

\end{tikzpicture}%
		\caption{Kullback-Leibler divergence according to \cref{eq:KL divergence} with $m=2000$.}
		\label{fig:kl sample based}
	\end{figure}
	
	The result depends on the choice of $m$ and on the feature that is used. \Cref{tab:wang result} lists the results for different values of $m$, where both the speed and the acceleration are considered as features. Considering the estimations regarding the different features, Wang \etal\ use the maximum as final estimation of the appropriate amount of data. The maximum required time is shown in the last column of \cref{tab:wang result}. \Cref{tab:wang result} shows that the final estimated of the required data heavily depends on $m$.
		
	\begin{table}
		\centering
		\caption{Appropriate data according to \cref{eq:completeness wang} with $\epsilon=10^{-4}$.}
		\label{tab:wang result}
		\begin{tabular}{rrrr}
			\toprule
			& \multicolumn{2}{c}{Feature} & \\
			$m$ & Speed & Acceleration & Maximum \\ \otoprule
			500 & $\unit[7]{min}$ & $\unit[5]{min}$ & $\unit[7]{min}$ \\
			1000 & $\unit[7]{min}$ & $\unit[12]{min}$ & $\unit[12]{min}$ \\
			2000 & $\unit[27]{min}$ & $\unit[23]{min}$ & $\unit[27]{min}$ \\
			4000 & $\unit[67]{min}$ & $\unit[27]{min}$ & $\unit[67]{min}$ \\
			6000 & $\unit[50]{min}$ & $\unit[40]{min}$ & $\unit[50]{min}$ \\
			8000 & $\unit[93]{min}$ & $\unit[40]{min}$ & $\unit[93]{min}$ \\
			10000 & $\unit[100]{min}$ & $\unit[33]{min}$ & $\unit[100]{min}$ \\
			\bottomrule
		\end{tabular}
	\end{table}
\end{example}


\subsection{Likelihood of new samples}
\label{sec:likelihood sample based}

This method is similar to the previous method in that a PDF is estimated using $n$ samples, i.e., $\hat{f}(\cdot;n)$. But instead of comparing $\hat{f}(\cdot;n)$ with $\hat{f}(\cdot;n+m)$, the $m$ new samples are directly compared with first $n$ samples by determining the likelihood of the new $m$ samples based on the estimated PDF $\hat{f}(\cdot;n)$. 

Assuming that the samples are independent, the likelihood is simply the product of $\hat{f}(x_i;n)$ with $i\in \{n+1, \ldots, n+m\}$. To get a value that is less dependent on the choice of $m$, the following value is used:
\begin{equation} \label{eq:likelihood estimation}
	J(n) = \left( \prod_{i=n+1}^{n+m} \hat{f}(x_i;n) \right)^{\frac{1}{m}}.
\end{equation}

Let us consider the log of \cref{eq:likelihood estimation}:
\begin{equation} \label{eq:loglikelihood estimation}
	\ln J(n) = \frac{1}{m} \sum_{i=n+1}^{n+m} \ln \hat{f}(x_i;n).
\end{equation}
An important choice is when $J(n)$ is good enough. The goal is to estimate the underlying distribution $f(x)$ as accurately as possible. Therefore, when the true distribution $f(x)$ is used instead of $\hat{f}(x_i;n)$ in \cref{eq:loglikelihood estimation}, the expected value is
\begin{equation} \label{eq:minus entropy}
	\ln J^* = \expectation{ \ln f(x) }.
\end{equation}
Notice that \cref{eq:minus entropy} equals the negative entropy of the $f(x)$. The entropy $\entropy{x}$ is defined as follows \cite{bishop2006pattern}:
\begin{equation} \label{eq:entropy}
	\entropy{x} = -\int f(x) \ln f(x) \ud x,
\end{equation}
where use is made of $f(x) \ln f(x) = 0$ if $f(x)=0$. Since the distribution $f(x)$ is unknown, the entropy can be estimated using $\hat{f}(x_i;n)$:
\begin{equation} \label{eq:entropy estimation}
	\hatentropy{x; n} = -\int \hat{f}(x;n) \ln \hat{f}(x;n) \ud x.
\end{equation}
The estimated entropy \cref{eq:entropy estimation} can be used to determine the desired value of the log-likelihood \eqref{eq:loglikelihood estimation}.

\begin{remark}
	Using the estimated entropy \eqref{eq:entropy estimation} promotes the use of high values of the bandwidth $h$, as this will result in higher values of the entropy. As a result, even though the estimated PDF $\hat{f}(x_i;n)$ might be worse with a higher value of $h$, the log-likelihood \cref{eq:loglikelihood estimation} might be closer to the estimated entropy \eqref{eq:entropy estimation}.
	
	A work-around could be to determine the log-likelihood of the first $n$ samples, i.e.,
	\begin{equation}
		\frac{1}{n} \sum_{i=1}^n \ln \hat{f} (x_i; n).
	\end{equation}
	If this value is significantly higher than $-\hatentropy{x;n}$, this might be an indication that $\hat{f}(x_i; n)$ is an inaccurate estimation of $f(x)$.
\end{remark}

\begin{example} \label{example:likelihood sample based}
	Based on the artificial data shown in \cref{fig:artificial data}, the log-likelihood \eqref{eq:loglikelihood estimation} is computed with $m=2000$, see \cref{fig:likelihood sample based}. It can be seen that the log-likelihood $\ln J(n)$ is on average the same as the negative entropy $-\hatentropy{x;n}$. However, there is not a clear trend visible regarding the log-likelihood and the amount of data.
	
	\setlength\figurewidth{0.6\linewidth}
	\setlength\figureheight{0.7\figurewidth}
	\begin{figure}
		\centering
		% This file was created by matplotlib2tikz v0.6.16.
\begin{tikzpicture}

\begin{axis}[
xlabel={Time [min]},
xmin=-10.5, xmax=293.833333333333,
ymin=-3, ymax=0,
width=\figurewidth,
height=\figureheight,
tick align=outside,
tick pos=left,
xmajorgrids,
x grid style={white!69.01960784313725!black},
ymajorgrids,
y grid style={white!69.01960784313725!black},
legend style={at={(0.97,0.03)}, anchor=south east, draw=white!80.0!black},
legend cell align={left},
legend entries={{$\ln J(n)$},{$-\hat{\textup{H}}[x;n]$}}
]
\addlegendimage{no markers, blue}
\addlegendimage{no markers, red}
\addplot [semithick, blue]
table {%
3.33333333333333 -1.48286248277346
6.66666666666667 -0.990783567932304
10 -2.1473927421105
13.3333333333333 -7.21906738478802
16.6666666666667 -0.999290209681654
20 -0.976406204128857
23.3333333333333 -0.896949033875354
26.6666666666667 -0.635435889892345
30 -1.55749449179073
33.3333333333333 -0.993774721306692
36.6666666666667 -1.3920493150687
40 -2.70763089229218
43.3333333333333 -2.17900679045107
46.6666666666667 -1.14942585188072
50 -0.445183452585065
53.3333333333333 -1.00443832001605
56.6666666666667 -1.1307727061787
60 -0.907587222895733
63.3333333333333 -0.850806381437652
66.6666666666667 -0.660235667701102
70 -0.879379966086497
73.3333333333333 -1.36129153557706
76.6666666666667 -1.52363730051123
80 -1.93541400711929
83.3333333333333 -1.60503103519867
86.6666666666667 -1.86392110903722
90 -0.946735730615624
93.3333333333333 -1.68141349276962
96.6666666666667 -1.16968814256937
100 -1.81631300005136
103.333333333333 -0.6325818472057
106.666666666667 -0.558225145976565
110 -0.538029203664425
113.333333333333 -0.7772045125717
116.666666666667 -1.71689485680577
120 -1.1951851708385
123.333333333333 -1.13997087896382
126.666666666667 -0.841326164184628
130 -1.75971304936648
133.333333333333 -0.734224039821198
136.666666666667 -1.07846976275205
140 -1.52350203148229
143.333333333333 -0.830732132464435
146.666666666667 -1.07921950874641
150 -1.73511322245676
153.333333333333 -1.67745678742317
156.666666666667 -1.9981579663809
160 -1.68027743294584
163.333333333333 -0.822651414430358
166.666666666667 -0.514967682265838
170 -0.488002552976629
173.333333333333 -0.491005334336116
176.666666666667 -1.06718473291285
180 -1.50277574858129
183.333333333333 -1.02404523610577
186.666666666667 -1.29825701781969
190 -0.939550200033033
193.333333333333 -1.68797190384103
196.666666666667 -1.62946099951827
200 -1.4612158178016
203.333333333333 -0.525981744544286
206.666666666667 -1.42295443642174
210 -1.15990181556058
213.333333333333 -0.941920155459013
216.666666666667 -1.24271823706866
220 -1.57439728776353
223.333333333333 -1.18362039331409
226.666666666667 -0.701115715151873
230 -1.27459608959546
233.333333333333 -1.0000306173186
236.666666666667 -1.14520121974918
240 -0.473800082467148
243.333333333333 -1.16587191443865
246.666666666667 -1.61047954032941
250 -1.55193913896079
253.333333333333 -1.51549847702367
256.666666666667 -1.29236228785889
260 -0.729455814614022
263.333333333333 -1.00051541792664
266.666666666667 -0.953335324593665
270 -1.39473791172298
273.333333333333 -1.05370830393162
276.666666666667 -1.6380693848862
280 -1.61670540451658
};
\addplot [semithick, red, dashed]
table {%
3.33333333333333 -1.18194051230298
6.66666666666667 -1.21513605970132
10 -1.19019431810883
13.3333333333333 -1.21640533709966
16.6666666666667 -1.26077591413288
20 -1.23760609898756
23.3333333333333 -1.22579887116951
26.6666666666667 -1.20000777156472
30 -1.15276298501137
33.3333333333333 -1.17915273522609
36.6666666666667 -1.17336005922424
40 -1.18320723543352
43.3333333333333 -1.20002306758842
46.6666666666667 -1.20668612367673
50 -1.21848812282469
53.3333333333333 -1.17967636101798
56.6666666666667 -1.18430609585409
60 -1.19534294835365
63.3333333333333 -1.1905797814248
66.6666666666667 -1.18383432537729
70 -1.16830783863315
73.3333333333333 -1.1634342052129
76.6666666666667 -1.17600338355551
80 -1.1780173714702
83.3333333333333 -1.18510152858963
86.6666666666667 -1.20161062166379
90 -1.21133461949712
93.3333333333333 -1.20498544390289
96.6666666666667 -1.20874468240174
100 -1.21447473173206
103.333333333333 -1.21159836231946
106.666666666667 -1.19792751348103
110 -1.18143603745015
113.333333333333 -1.16794348570582
116.666666666667 -1.16408379221875
120 -1.16492404987804
123.333333333333 -1.15963085580799
126.666666666667 -1.16667786044134
130 -1.16286624654089
133.333333333333 -1.16129678815331
136.666666666667 -1.15135728831149
140 -1.15692337164738
143.333333333333 -1.16924280040115
146.666666666667 -1.16658213151312
150 -1.16839996200871
153.333333333333 -1.17687863945582
156.666666666667 -1.17656702451406
160 -1.17269337959715
163.333333333333 -1.16926522300236
166.666666666667 -1.16572232241086
170 -1.15603001404736
173.333333333333 -1.14544281238822
176.666666666667 -1.13610640899826
180 -1.13991800695481
183.333333333333 -1.14609815127386
186.666666666667 -1.14835780521989
190 -1.15657359840536
193.333333333333 -1.15741575310426
196.666666666667 -1.15647115226701
200 -1.15739305633158
203.333333333333 -1.16171158488703
206.666666666667 -1.15380273039097
210 -1.15778376451625
213.333333333333 -1.1553490361616
216.666666666667 -1.15517605559039
220 -1.16002481460394
223.333333333333 -1.16619337978566
226.666666666667 -1.17042651488939
230 -1.16598214016866
233.333333333333 -1.16620596959223
236.666666666667 -1.16638200330243
240 -1.16991707021804
243.333333333333 -1.16201366962919
246.666666666667 -1.16391821719306
250 -1.16070469993504
253.333333333333 -1.1643139893881
256.666666666667 -1.16781750722612
260 -1.16639695338033
263.333333333333 -1.16295682396678
266.666666666667 -1.16387791022929
270 -1.16160594458116
273.333333333333 -1.1638388847993
276.666666666667 -1.16480318712557
280 -1.16480795745098
};
\end{axis}

\end{tikzpicture}%
		\caption{Log-likelihood \cref{eq:loglikelihood estimation} with $m=2000$ and the estimated entropy \cref{eq:entropy estimation}.}
		\label{fig:likelihood sample based}
	\end{figure}
\end{example}


\section{Activity-based methods}
\label{sec:activity based methods}


\subsection{Comparing probability density functions}
\label{sec:comparing pdf activities}

Instead of considering the PDFs of the time samples (see \cref{sec:sample based pdf}), the PDFs of the activities' features will be considered in this section. An activity refers to the behavior of the vehicle between two events. As the number of activities is much less than the number of samples, the analysis using the activities is much less computationally intensive.

For illustrating the methods, three features that describe an activity are considered:
\begin{itemize}
	\item $v_{\textup{end}}$: The end speed of the activity.
	\item $\Delta v$: The difference between the end and start speed.
	\item $\Delta t$: The duration of the activity.
\end{itemize}

For each feature, a univariate PDF is estimated using Kernel Density Estimation (KDE), see \cref{eq:kde}, with the bandwidth according to Silverman's rule, see \cref{eq:silverman}. For each feature, the appropriate number of activities is estimated using the Kullback-Leibler divergence, according to \cref{eq:kl epsilon}, resulting in $n_{v_{\textup{end}}}$, $n_{\Delta v}$, and $n_{\Delta t}$. The total number of activities is simply the maximum of $n_{v_{\textup{end}}}$, $n_{\Delta v}$, and $n_{\Delta t}$.

\begin{example}
	Based on the artificial data, the KL divergence \cref{eq:KL divergence} is computed. \Cref{fig:kl activity based} shows the result for $m=10$. Using $\epsilon = 10^{-3}$ results in $n_{v_{\textup{end}}}=150$, $n_{\Delta v}=80$, and $n_{\Delta t}=20$ which correspond to $\unit[46]{min}$, $\unit[25]{min}$, and $\unit[7]{min}$, respectively. Therefore, according to this method, the appropriate amount of data is $\unit[46]{min}$. 
	
	\setlength\figurewidth{0.6\linewidth}
	\setlength\figureheight{0.7\figurewidth}
	\begin{figure}
		\centering
		% This file was created by matplotlib2tikz v0.6.14.
\begin{tikzpicture}

\begin{axis}[
xlabel={Time [min]},
ylabel={KL divergence},
xmin=-10.6114166666667, xmax=293.826416666667,
ymin=5.506419140782e-07, ymax=0.645093865006154,
ymode=log,
width=\figurewidth,
height=\figureheight,
tick align=outside,
tick pos=left,
xmajorgrids,
x grid style={lightgray!92.026143790849673!black},
ymajorgrids,
y grid style={lightgray!92.026143790849673!black},
legend entries={{$v_{\textup{end}}$},{$\Delta v$},{$\Delta t$}},
legend cell align={left},
legend style={draw=white!80.0!black}
]
\addlegendimage{no markers, blue}
\addlegendimage{no markers, red}
\addlegendimage{no markers, green!50.0!black}
\addplot [semithick, blue]
table {%
3.22666666666667 0.0542095584884025
7.03166666666667 0.0269425177766985
9.51333333333333 0.0307954541956445
13.2833333333333 0.200111536439902
15.2183333333333 0.00150628676752472
19.3116666666667 0.0156757584789558
20.955 0.00251683366301053
25.4433333333333 0.00547169774010929
28.875 0.00204410639970234
32.4783333333333 0.00291486948558313
35.8016666666667 0.00329098959411608
38.555 0.0063873713641226
41.2883333333333 0.0188105046758615
43.7 0.00720693933817983
45.8583333333333 0.00437119933982359
48.9416666666667 0.000136046066171519
52.62 0.00243119155686667
56.0516666666667 0.00200002922739795
58.6816666666667 0.00286163306354062
61.9083333333333 0.00158610402641542
64.215 0.00125018936247013
66.5166666666667 0.00121544797903994
69.2616666666667 0.000532318365874471
74.175 0.00248898493452929
76.2166666666667 0.000120351943888532
80.58 0.00108454546264998
82.835 0.00118663827579269
86.845 0.00147120015505199
89.09 -0.000336962079867676
93.72 0.000195376933542005
97.545 0.0012747369237543
100.02 0.00111798591925878
102.33 0.000171810395664846
105.761666666667 0.00010973752423784
109.115 0.000469582925447773
112.37 9.24283784526127e-05
115.588333333333 0.000429305670348922
119.293333333333 0.000788485723688211
121.851666666667 0.000612544395229938
124.501666666667 0.000626083441576879
127.805 0.000369112414945815
130.513333333333 0.000625518833292556
132.728333333333 -0.000121138453322974
136.228333333333 0.000549218947048726
138.8 0.000972587180079378
140.811666666667 0.000883196230468117
142.7 0.000442636537358956
145.056666666667 0.000433640408271801
147.041666666667 0.000306194786321953
149.461666666667 0.00069000381673346
151.1 0.000574949980978689
152.615 0.000590948435104172
154.468333333333 0.000409064752186542
156.353333333333 0.000349221894730092
158.696666666667 0.000361373693316314
160.356666666667 0.000255854469144491
162.825 0.000234978723233373
165.365 -1.77931143684999e-05
168.71 -3.30812345192335e-05
172.803333333333 -1.16961563864695e-05
176.4 0.00028180452096596
178.83 0.000299964751093749
182.02 0.000388319213995244
184.616666666667 0.000336317978226319
186.551666666667 0.000432142404199282
189.38 0.000264928110645024
192.69 0.000144606083971795
195.335 0.000199357062273433
198.151666666667 0.000289507390278612
201.988333333333 5.9506652877414e-05
205.976666666667 0.000329493541665999
208.3 0.00015886432117747
210.438333333333 -3.68409959864171e-05
215.353333333333 4.68240805110945e-05
219.745 0.000363170008322476
222.001666666667 0.000136863488685619
225.611666666667 0.000322799785110594
228.073333333333 1.03925399971253e-06
230.448333333333 0.000224524054187803
232.358333333333 -1.66903404311134e-05
236.365 0.000162149365825985
239.905 -5.45588837400854e-06
242.72 1.20135990238428e-05
247.211666666667 0.000146668894555989
249.07 0.000192765884365893
250.603333333333 0.000282717917996984
252.301666666667 0.000129764961029674
254.623333333333 0.000232098661813182
256.77 -4.67662468583281e-05
259.576666666667 5.65323548360028e-05
262.866666666667 5.1993361346559e-05
265.441666666667 4.37514566804885e-05
269.91 0.000259005371937934
272.08 -7.07367845812407e-05
275.341666666667 7.89316540681932e-05
277.97 9.22405768953312e-05
279.988333333333 8.38165025983096e-05
};
\addplot [semithick, red, dashed]
table {%
3.22666666666667 0.105370793440414
7.03166666666667 0.0146445730283461
9.51333333333333 0.00305670246743903
13.2833333333333 0.00597601279861465
15.2183333333333 0.236925939249084
19.3116666666667 0.0024188880495738
20.955 0.00472052425521063
25.4433333333333 0.00168850932337503
28.875 0.000889030082198155
32.4783333333333 0.0108150852877953
35.8016666666667 0.00122999406547993
38.555 0.000728917213305437
41.2883333333333 0.000440140132318052
43.7 0.000519026757654447
45.8583333333333 0.000469052351383896
48.9416666666667 0.00228252127500684
52.62 0.000448902805713486
56.0516666666667 0.00074049380386259
58.6816666666667 0.000396967378371502
61.9083333333333 0.00033496293449209
64.215 0.000946444248279994
66.5166666666667 0.000230437836737747
69.2616666666667 0.000419527854537567
74.175 0.000221395573939673
76.2166666666667 0.0014859277113139
80.58 0.00023713495228256
82.835 0.000266536639856686
86.845 0.000209919582824012
89.09 0.00193198916470405
93.72 0.000624273201891977
97.545 0.000153227945881965
100.02 0.000989043184069464
102.33 0.000332136258949337
105.761666666667 0.000146236103877856
109.115 0.000117842376111587
112.37 0.00012160809983456
115.588333333333 0.000591560403886081
119.293333333333 9.30894968217924e-05
121.851666666667 0.000206739734264301
124.501666666667 0.000160675620774985
127.805 8.82043501439129e-05
130.513333333333 8.00265145454859e-05
132.728333333333 0.000165780060394222
136.228333333333 9.84057935686925e-05
138.8 7.98794855311743e-05
140.811666666667 0.000100025418719055
142.7 0.000115408909123093
145.056666666667 4.83961701731399e-05
147.041666666667 4.70440574229964e-05
149.461666666667 6.0425576991385e-05
151.1 6.89099501955175e-05
152.615 6.72070818228976e-05
154.468333333333 4.36881967619196e-05
156.353333333333 4.22640311540958e-05
158.696666666667 5.46015932429366e-05
160.356666666667 3.72673875137014e-05
162.825 0.000215309523826621
165.365 3.37889883556702e-05
168.71 0.00011511239665379
172.803333333333 2.99744874555962e-05
176.4 5.08381205515447e-05
178.83 4.00631088638104e-05
182.02 7.59666489871828e-05
184.616666666667 4.40785947721533e-05
186.551666666667 7.74344973032487e-05
189.38 0.000137886141164909
192.69 4.14029637201341e-05
195.335 3.20274663713743e-05
198.151666666667 0.000165403650640751
201.988333333333 0.000350669610407988
205.976666666667 5.00077366329718e-05
208.3 2.41645615629222e-05
210.438333333333 0.000228066989454603
215.353333333333 8.63412129577112e-05
219.745 1.99512682334981e-05
222.001666666667 0.000246848888247636
225.611666666667 2.98116237986765e-05
228.073333333333 2.7819161395854e-05
230.448333333333 3.63396247767908e-05
232.358333333333 0.000193804747304565
236.365 2.32719819107043e-05
239.905 3.84858350968604e-05
242.72 6.02228080549367e-05
247.211666666667 1.50634487086674e-05
249.07 2.30344994892708e-05
250.603333333333 1.61120090384485e-05
252.301666666667 3.640067902219e-05
254.623333333333 3.0645770023208e-05
256.77 0.000177668781699228
259.576666666667 1.81785963182981e-05
262.866666666667 2.17625446680808e-05
265.441666666667 0.000119110221195449
269.91 2.52427614160914e-05
272.08 6.32329398678457e-05
275.341666666667 0.000127630279269239
277.97 1.75983460846102e-05
279.988333333333 1.60996323219756e-05
};
\addplot [semithick, green!50.0!black, dash pattern=on 1pt off 3pt on 3pt off 3pt]
table {%
3.22666666666667 0.341798752456425
7.03166666666667 0.0179259331646596
9.51333333333333 0.000704769656509018
13.2833333333333 0.0116598777079103
15.2183333333333 0.00858342801477595
19.3116666666667 0.0105507487162978
20.955 0.00470144063235658
25.4433333333333 0.00949787482640373
28.875 0.00319274355020972
32.4783333333333 0.000977978876592609
35.8016666666667 0.00161740253509904
38.555 0.00245676157648896
41.2883333333333 0.00251125865392462
43.7 0.00172355344295518
45.8583333333333 0.000379555900715479
48.9416666666667 0.000805219086202567
52.62 0.000856857314998332
56.0516666666667 0.000400964056872536
58.6816666666667 0.000560177651910222
61.9083333333333 0.000490555006853933
64.215 0.000936703824539929
66.5166666666667 0.000671099522402341
69.2616666666667 0.272743974927835
74.175 0.00102268579548896
76.2166666666667 0.00160955745913207
80.58 0.000432119364174067
82.835 0.0133954094697628
86.845 0.000429245476178375
89.09 0.00108906812717302
93.72 0.00650607011692267
97.545 0.000764396137152473
100.02 0.00046453020355274
102.33 0.000198372561465508
105.761666666667 0.000552383324320021
109.115 0.000424685916812687
112.37 0.000464374010981089
115.588333333333 0.000540223672378426
119.293333333333 0.000353157039979881
121.851666666667 0.000316689990141232
124.501666666667 0.000341734905035272
127.805 0.000341107864768241
130.513333333333 0.000232012663109054
132.728333333333 0.000132754723114557
136.228333333333 0.000267244814162629
138.8 0.000285794952340916
140.811666666667 0.000206015391289059
142.7 0.000215550399080574
145.056666666667 0.000331402775817406
147.041666666667 0.00019844233870364
149.461666666667 0.000232048807372125
151.1 0.000186799097557377
152.615 0.000231570527917557
154.468333333333 0.000216491870503763
156.353333333333 0.000175814844186111
158.696666666667 0.000201720756824708
160.356666666667 0.000154954414072725
162.825 0.000162343317294009
165.365 9.53168953870376e-05
168.71 0.000128807773558996
172.803333333333 7.24367666032527e-05
176.4 8.68289024176973e-05
178.83 0.000128082273039187
182.02 0.000108196917533899
184.616666666667 8.15854745705458e-05
186.551666666667 6.79454896389963e-05
189.38 6.8258064974945e-05
192.69 6.07073567559154e-05
195.335 7.7746465687508e-05
198.151666666667 0.000251700395772957
201.988333333333 0.000125424097351736
205.976666666667 7.41595544024375e-05
208.3 5.10198405960503e-05
210.438333333333 0.000109461086854468
215.353333333333 0.000777924425813871
219.745 7.86799193115315e-05
222.001666666667 8.66148171026227e-05
225.611666666667 8.47037601848355e-05
228.073333333333 2.49507027104002e-05
230.448333333333 6.54492658022858e-05
232.358333333333 0.000292057123648238
236.365 0.000112534110369363
239.905 7.60504081648544e-05
242.72 0.000106669564646298
247.211666666667 8.20602622585541e-05
249.07 8.11961151299545e-05
250.603333333333 6.87160019523139e-05
252.301666666667 6.72019041409861e-05
254.623333333333 4.41176045383966e-05
256.77 6.13521753887018e-05
259.576666666667 8.3297855216401e-05
262.866666666667 6.44407050871917e-05
265.441666666667 7.15923388529513e-05
269.91 4.88966702249473e-05
272.08 4.65358860283515e-05
275.341666666667 4.21786361112641e-05
277.97 3.74686229533874e-05
279.988333333333 5.16950452266526e-05
};
\end{axis}

\end{tikzpicture}%
		\caption{Kullback-Leibler divergence \cref{eq:KL divergence} using the features from the activities with $m=10$.}
		\label{fig:kl activity based}
	\end{figure}

	The result depends hugely on the choice of $m$. Obviously, when more data is added, the difference between the two PDFs $\hat{f}(x;n+m)$ and $\hat{f}(x;n)$ is likely to be higher. For example, when $m=5$ is used, the estimated amount of appropriate data is $\unit[21]{min}$ instead of $\unit[46]{min}$. When $m=15$, this is $\unit[74]{min}$.
\end{example}


\subsection{Likelihood of new activities}
\label{sec:likelihood new activities}

This method is similar to the method presented in \cref{sec:likelihood sample based}, but instead of computing the likelihood of the samples, the likelihood of the activities is computed using the features of the activities. 

\begin{example}
	Based on the artificial data, the likelihood of the feature $\Delta v$ is computed, see \cref{fig:likelihood activity based}. Similarly when applying this method for the samples (see \cref{example:likelihood sample based}), there is not a clear trend visible regarding the log-likelihood and the amount of data.
	
	\setlength\figurewidth{0.6\linewidth}
	\setlength\figureheight{0.7\figurewidth}
	\begin{figure}
		\centering
		% This file was created by matplotlib2tikz v0.6.14.
\begin{tikzpicture}

\begin{axis}[
xlabel={Time [min]},
xmin=-10.71025, xmax=295.901916666667,
ymin=-3.24098845099254, ymax=-0.54089124750259,
width=\figurewidth,
height=\figureheight,
tick align=outside,
tick pos=left,
xmajorgrids,
x grid style={lightgray!92.026143790849673!black},
ymajorgrids,
y grid style={lightgray!92.026143790849673!black},
legend entries={{$\ln J(n)$},{$-\hat{\textup{H}}[x;n]$}},
legend style={at={(0.97,0.03)}, anchor=south east, draw=white!80.0!black},
legend cell align={left}
]
\addlegendimage{no markers, blue}
\addlegendimage{no markers, red}
\addplot [semithick, blue]
table {%
3.22666666666667 -1.68438713847269
7.03166666666667 -1.20908087916945
9.51333333333333 -1.25577235376987
13.2833333333333 -1.24229079995518
15.2183333333333 -3.11825675992482
19.3116666666667 -1.01484165746668
20.955 -1.63638605773551
25.4433333333333 -1.00434817415724
28.875 -1.12272098849379
32.4783333333333 -1.85986362308317
35.8016666666667 -1.01247358273794
38.555 -0.998970180368461
41.2883333333333 -1.05527439585289
43.7 -0.956146672586978
45.8583333333333 -1.23949229480806
48.9416666666667 -1.62317550588898
52.62 -0.95122470830587
56.0516666666667 -1.48495711126698
58.6816666666667 -1.03255759420302
61.9083333333333 -0.865839379658711
64.215 -1.47264813899929
66.5166666666667 -0.952652562202561
69.2616666666667 -1.32220707519382
74.175 -0.913126614046882
76.2166666666667 -1.62414632562245
80.58 -1.34349245626765
82.835 -1.3572680214183
86.845 -1.22732336969712
89.09 -1.7766527037161
93.72 -1.62525711520725
97.545 -1.26745721285821
100.02 -1.34252851088723
102.33 -1.41030642730611
105.761666666667 -0.831135539004624
109.115 -0.889785127951881
112.37 -1.19737864897683
115.588333333333 -1.63580539985621
119.293333333333 -0.913395877709148
121.851666666667 -1.38658434777234
124.501666666667 -1.12126504299546
127.805 -1.05104578983041
130.513333333333 -0.984181880043323
132.728333333333 -1.12507753558692
136.228333333333 -0.836346844045767
138.8 -0.995052127345691
140.811666666667 -0.787001474621974
142.7 -0.710102365113068
145.056666666667 -1.06990274719241
147.041666666667 -0.949847921006493
149.461666666667 -0.93458918120409
151.1 -0.827297161772912
152.615 -0.790040352738144
154.468333333333 -0.884729301172547
156.353333333333 -0.90762281648812
158.696666666667 -0.779797881578626
160.356666666667 -0.993908473111251
162.825 -1.13390597728515
165.365 -0.997810763930262
168.71 -1.45935491265409
172.803333333333 -0.930252808337346
176.4 -1.06874176335649
178.83 -1.1826531501231
182.02 -0.7200991664831
184.616666666667 -0.961288023256119
186.551666666667 -1.30953107774663
189.38 -1.32472811941393
192.69 -1.20431018198292
195.335 -1.13417491188099
198.151666666667 -1.53238348657598
201.988333333333 -1.28161615559442
205.976666666667 -1.14015706723629
208.3 -1.06396119025699
210.438333333333 -1.33493198861482
215.353333333333 -1.43180822339282
219.745 -1.00113725054542
222.001666666667 -1.63806490192862
225.611666666667 -0.766026053002339
228.073333333333 -0.800523678636485
230.448333333333 -1.15940699550164
232.358333333333 -1.61915301300749
236.365 -1.06489127596154
239.905 -1.28245706917963
242.72 -1.20773767274121
247.211666666667 -1.01738411021235
249.07 -0.848502310147996
250.603333333333 -0.833620799651584
252.301666666667 -1.26773699064629
254.623333333333 -1.11583742138364
256.77 -1.41368053083015
259.576666666667 -1.08783822940592
262.866666666667 -1.22166384552343
265.441666666667 -1.68316694479935
269.91 -1.14721301927089
272.08 -1.30738675878648
275.341666666667 -1.57017575724485
277.97 -0.909057140395858
279.988333333333 -0.803869654448426
281.965 -0.663622938570315
};
\addplot [semithick, red, dashed]
table {%
3.22666666666667 -1.34207578005886
7.03166666666667 -1.39710615988595
9.51333333333333 -1.31389958424723
13.2833333333333 -1.31683238184792
15.2183333333333 -1.31554374836667
19.3116666666667 -1.2094949158551
20.955 -1.2096031908715
25.4433333333333 -1.23015959717421
28.875 -1.22701458394564
32.4783333333333 -1.21244874894147
35.8016666666667 -1.20375881544277
38.555 -1.18786956595273
41.2883333333333 -1.17499981200367
43.7 -1.1723725006487
45.8583333333333 -1.16713771750985
48.9416666666667 -1.17695313426939
52.62 -1.18343424819649
56.0516666666667 -1.17730686949495
58.6816666666667 -1.19053557063296
61.9083333333333 -1.18913253609767
64.215 -1.18220638614203
66.5166666666667 -1.19067193916242
69.2616666666667 -1.18712689466111
74.175 -1.18302896014843
76.2166666666667 -1.17860152907832
80.58 -1.16113705944544
82.835 -1.16251180430426
86.845 -1.1719366706028
89.09 -1.17655452258157
93.72 -1.16451003082551
97.545 -1.15320748822842
100.02 -1.155050551521
102.33 -1.14931514327761
105.761666666667 -1.13834908190287
109.115 -1.13484706239756
112.37 -1.13305031921985
115.588333333333 -1.13663949030852
119.293333333333 -1.136214993121
121.851666666667 -1.13516393372282
124.501666666667 -1.12407309742805
127.805 -1.1267877695153
130.513333333333 -1.12531519581393
132.728333333333 -1.12462034025968
136.228333333333 -1.11165580472933
138.8 -1.11024300083508
140.811666666667 -1.11123592457318
142.7 -1.10918610988983
145.056666666667 -1.10595985929141
147.041666666667 -1.10650677539651
149.461666666667 -1.1073545354179
151.1 -1.10669025125246
152.615 -1.10530549967375
154.468333333333 -1.10315759491945
156.353333333333 -1.10311951359023
158.696666666667 -1.10301717028257
160.356666666667 -1.10123275085419
162.825 -1.10287802183698
165.365 -1.10472552902216
168.71 -1.10601332172665
172.803333333333 -1.11274930311441
176.4 -1.11250924906637
178.83 -1.11356986897581
182.02 -1.11438928564389
184.616666666667 -1.11220459490623
186.551666666667 -1.1129938248621
189.38 -1.11282725715547
192.69 -1.11679396041125
195.335 -1.1155228507272
198.151666666667 -1.11796844768381
201.988333333333 -1.11797972229757
205.976666666667 -1.11072064612373
208.3 -1.11232943984291
210.438333333333 -1.11340925160902
215.353333333333 -1.10663193825388
219.745 -1.10733339980658
222.001666666667 -1.1079820175158
225.611666666667 -1.10266313552967
228.073333333333 -1.10127770089169
230.448333333333 -1.10029385867881
232.358333333333 -1.09976854243314
236.365 -1.10179126459382
239.905 -1.10248651056162
242.72 -1.10571880387525
247.211666666667 -1.10312254078897
249.07 -1.10386174409229
250.603333333333 -1.10304782855867
252.301666666667 -1.10239976051386
254.623333333333 -1.10533285222175
256.77 -1.1070342716456
259.576666666667 -1.09932397919869
262.866666666667 -1.10078160761242
265.441666666667 -1.10327448026145
269.91 -1.09832925923436
272.08 -1.09967981579554
275.341666666667 -1.093651651053
277.97 -1.09102777692843
279.988333333333 -1.09076159972411
281.965 -1.09007003853569
};
\end{axis}

\end{tikzpicture}%
		\caption{Log-likelihood \cref{eq:loglikelihood estimation} with $m=10$ and the estimated entropy \cref{eq:entropy estimation} using the feature $\Delta v$.}
		\label{fig:likelihood activity based}
	\end{figure}

	A possible reason why the result in \cref{fig:likelihood activity based} is inconclusive could be the fact that the chosen bandwidth $h$ is too large and therefore the estimated PDF $\hat{f}(x;n)$ is oversmoothed. When using pseudo likelihood cross-validation \cite{turlach1993bandwidthselection} to determine the optimal bandwidth, a smaller bandwidth is obtained, see \cref{fig:bandwidth}. This results in a different log-likelihood, see \cref{fig:loglikelihood activities small bandwidth}.
	
	\setlength\figureheight{0.4\linewidth}
	\setlength\figurewidth{0.48\linewidth}
	\begin{figure}
		\centering
		\begin{minipage}{0.48\linewidth}
			\centering
			% This file was created by matplotlib2tikz v0.6.14.
\begin{tikzpicture}

\begin{axis}[
xlabel={Time [min]},
xmin=-10.71025, xmax=295.901916666667,
ymin=-2.67294835928751, ymax=0.680650299092094,
width=\figurewidth,
height=\figureheight,
tick align=outside,
tick pos=left,
xmajorgrids,
x grid style={lightgray!92.026143790849673!black},
ymajorgrids,
y grid style={lightgray!92.026143790849673!black},
legend style={at={(0.03,0.97)}, anchor=north west, draw=white!80.0!black},
legend entries={{$\ln J(n)$},{$-\hat{\textup{H}}[x;n]$}},
legend cell align={left}
]
\addlegendimage{no markers, blue}
\addlegendimage{no markers, red}
\addplot [semithick, blue]
table {%
3.22666666666667 -1.6309234208491
7.03166666666667 -1.32998900150549
9.51333333333333 -1.28681404297673
13.2833333333333 -1.25676324480822
15.2183333333333 -2.48253119206962
19.3116666666667 -1.08986241381522
20.955 -1.6728414240503
25.4433333333333 -1.07787930244152
28.875 -1.20327887333041
32.4783333333333 -1.75538726561133
35.8016666666667 -1.01182866208123
38.555 -0.924219718569494
41.2883333333333 -0.980639425682897
43.7 -0.856171268018169
45.8583333333333 -1.13442018851451
48.9416666666667 -1.797915010171
52.62 -0.785280858504851
56.0516666666667 -1.44371552795523
58.6816666666667 -0.857990705540537
61.9083333333333 -0.571095670573732
64.215 -1.38943889053565
66.5166666666667 -0.59981839685359
69.2616666666667 -0.987523694820783
74.175 -0.710588701154528
76.2166666666667 -1.65166135844148
80.58 -0.827854453374787
82.835 -0.823236929215097
86.845 -0.735202505233458
89.09 -2.10875914576904
93.72 -1.29541180577215
97.545 -0.817985332362893
100.02 -2.52051205663389
102.33 -0.98825579555582
105.761666666667 -0.369394895130894
109.115 -0.254837133433848
112.37 -0.703022793140295
115.588333333333 -1.44381433678247
119.293333333333 -0.380144865073186
121.851666666667 -0.782480789293619
124.501666666667 -0.989016234518228
127.805 -0.355950451209122
130.513333333333 -0.63451519664867
132.728333333333 -0.48911439324261
136.228333333333 -0.0789439892110989
138.8 -0.388888154882807
140.811666666667 0.0151293257590765
142.7 0.0672981219608536
145.056666666667 -0.467362628333633
147.041666666667 -0.209054602199499
149.461666666667 -0.155358231406639
151.1 -0.104778211185639
152.615 -0.226219065320531
154.468333333333 -0.0513972008389445
156.353333333333 -0.249830957212506
158.696666666667 -0.0323259350044832
160.356666666667 -0.272898528004486
162.825 -0.459827955533993
165.365 -0.335411685355291
168.71 -1.18475907905305
172.803333333333 -0.372891112476155
176.4 -0.275651370111739
178.83 -0.663335319402489
182.02 -0.00531653996564696
184.616666666667 -0.29781065928378
186.551666666667 -0.743028222671254
189.38 -1.25128543342377
192.69 -0.408654100071241
195.335 -0.629915663888366
198.151666666667 -0.802545506351851
201.988333333333 -0.385462170565529
205.976666666667 -0.472405987856906
208.3 -0.547528726871033
210.438333333333 -1.9392757157569
215.353333333333 -0.561727440178073
219.745 -0.542543201035486
222.001666666667 -1.22474523128437
225.611666666667 0.00210825025077188
228.073333333333 -0.119773106718879
230.448333333333 -0.565404198633043
232.358333333333 -1.64958419478578
236.365 -0.110002809380603
239.905 -0.68468249036486
242.72 -0.232349694441279
247.211666666667 -0.468107033171013
249.07 0.010502377621721
250.603333333333 0.127213576732448
252.301666666667 -0.558249015171307
254.623333333333 -0.340882007738524
256.77 -0.649262942749809
259.576666666667 -0.228449628147758
262.866666666667 -0.696687501980657
265.441666666667 -1.0504912551558
269.91 -0.410972670521516
272.08 -0.436816381726453
275.341666666667 -0.837815736363292
277.97 0.0191764061137775
279.988333333333 0.239828638610478
281.965 0.528213996438475
};
\addplot [semithick, red, dashed]
table {%
3.22666666666667 -1.39401925210651
7.03166666666667 -1.46442804642201
9.51333333333333 -1.37751342054742
13.2833333333333 -1.37400905913389
15.2183333333333 -1.38206960885248
19.3116666666667 -1.31600419107752
20.955 -1.32158906079761
25.4433333333333 -1.31181269095255
28.875 -1.3064957158159
32.4783333333333 -1.2843837344512
35.8016666666667 -1.20325366054628
38.555 -1.12106452675535
41.2883333333333 -1.07155363182985
43.7 -1.06788233075532
45.8583333333333 -1.05335257418241
48.9416666666667 -1.04908230843971
52.62 -1.00037285564355
56.0516666666667 -0.942311587607039
58.6816666666667 -0.967677730824975
61.9083333333333 -0.963613403117718
64.215 -0.948749448765889
66.5166666666667 -0.972871647821039
69.2616666666667 -0.960835976478102
74.175 -0.947006857204914
76.2166666666667 -0.948968222832335
80.58 -0.823180576198896
82.835 -0.759306000006032
86.845 -0.767660719132667
89.09 -0.774669785223821
93.72 -0.757923753848693
97.545 -0.71291447442543
100.02 -0.719599935927831
102.33 -0.744801881869212
105.761666666667 -0.735132552403523
109.115 -0.71724455986951
112.37 -0.709831614474704
115.588333333333 -0.713726020798026
119.293333333333 -0.681649245924784
121.851666666667 -0.666617917969166
124.501666666667 -0.63970220767955
127.805 -0.658542002143111
130.513333333333 -0.656927613818354
132.728333333333 -0.660595012278177
136.228333333333 -0.642918537145517
138.8 -0.623442562422133
140.811666666667 -0.621910993416725
142.7 -0.613011006588821
145.056666666667 -0.629882714242598
147.041666666667 -0.603602581564317
149.461666666667 -0.60095809014942
151.1 -0.596843855763748
152.615 -0.579002424223009
154.468333333333 -0.577023159964598
156.353333333333 -0.571713961397016
158.696666666667 -0.569637384435422
160.356666666667 -0.590440424685952
162.825 -0.589158826437672
165.365 -0.577525584013789
168.71 -0.550797692048346
172.803333333333 -0.56225045769868
176.4 -0.563111621297098
178.83 -0.54979884565683
182.02 -0.665454066891655
184.616666666667 -0.660985354818554
186.551666666667 -0.546913687217271
189.38 -0.535494236401028
192.69 -0.566399717326383
195.335 -0.563689986868638
198.151666666667 -0.567777308342405
201.988333333333 -0.55517474609204
205.976666666667 -0.481064282812734
208.3 -0.47124502552621
210.438333333333 -0.487989770714476
215.353333333333 -0.481438674830345
219.745 -0.593737626003807
222.001666666667 -0.608714014305016
225.611666666667 -0.551944483145875
228.073333333333 -0.536503570076105
230.448333333333 -0.523721393631859
232.358333333333 -0.523371564276519
236.365 -0.387243822260011
239.905 -0.386742131204154
242.72 -0.428201494773515
247.211666666667 -0.425645673371382
249.07 -0.428725562473012
250.603333333333 -0.426631847543336
252.301666666667 -0.423399241232059
254.623333333333 -0.426748667197977
256.77 -0.415377953227028
259.576666666667 -0.345955740652277
262.866666666667 -0.347010750331692
265.441666666667 -0.376171618848423
269.91 -0.36318989995865
272.08 -0.475869422889173
275.341666666667 -0.495230962865836
277.97 -0.351389336240016
279.988333333333 -0.324608325877723
281.965 -0.309745528200771
};
\end{axis}

\end{tikzpicture}%
			\subcaption{Log-likelihood.}
			\label{fig:loglikelihood activities small bandwidth}
		\end{minipage}%
		\begin{minipage}{0.48\linewidth}
			\centering
			% This file was created by matplotlib2tikz v0.7.5.
\begin{tikzpicture}

\begin{axis}[
axis x line*=bottom,
axis y line*=left,
every x tick/.style={black},
every y tick/.style={black},
height=\figureheight,
log basis x={10},
log basis y={10},
tick align=outside,
tick pos=left,
width=\figurewidth,
x grid style={white!69.01960784313725!black},
xlabel={Number of samples},
xmin=100, xmax=5000,
xmode=log,
xtick style={color=black},
y grid style={white!69.01960784313725!black},
ylabel={Bandwidth},
ymin=0.12073463121632, ymax=0.514133563500785,
ymode=log,
ytick style={color=black}
]
\addplot [line width=0.7000000000000001pt, gray]
table {%
100 0.333902544590215
108 0.334746992843899
117 0.33249691495283
127 0.313299239565951
138 0.308128449857578
149 0.287087379333731
161 0.275852692097465
175 0.247590919809407
189 0.240371849561726
205 0.245019151799335
222 0.238213129893222
241 0.232594918722967
261 0.223326930772757
282 0.210746068090978
306 0.216497613329582
331 0.253957966161198
359 0.263488258335087
389 0.259249215169254
421 0.254883009145746
456 0.248740533199762
494 0.235012928468767
535 0.220218276842796
579 0.202869246927223
627 0.20290435756291
679 0.190874804682713
736 0.214007512145567
797 0.217916819764872
863 0.207377564314436
935 0.198103970787746
1013 0.195903264029502
1097 0.178583258596411
1188 0.182778529580626
1287 0.182014773896209
1394 0.185411253471078
1510 0.185224570599584
1635 0.173469219282776
1771 0.173703922947084
1918 0.172726937037639
2078 0.168995061714128
2250 0.161121478215138
2437 0.155267211210851
2640 0.16282926628745
2859 0.163813484923627
3097 0.158158303309362
3354 0.152430615691622
3633 0.144719178998213
3935 0.144723048908827
4262 0.155973985312413
4616 0.152543706717652
5000 0.148094988830411
};
\addplot [line width=0.7000000000000001pt, gray, dotted]
table {%
100 0.345479296425376
108 0.331495634249555
117 0.330685959154714
127 0.330651878039332
138 0.303374311500354
149 0.288045984767778
161 0.283579062143539
175 0.285759824891524
189 0.277264201851355
205 0.276537918621916
222 0.261552908075927
241 0.254902806261265
261 0.252412415046285
282 0.24735242620911
306 0.250473019101765
331 0.23635267972024
359 0.211317117665954
389 0.215603831320825
421 0.202279662613266
456 0.188193458503748
494 0.189605745892871
535 0.201713723998427
579 0.209045736912432
627 0.211462713508606
679 0.216255722295364
736 0.221059383003018
797 0.222508953506758
863 0.225867820591947
935 0.212422814712975
1013 0.211964288180541
1097 0.194652580532807
1188 0.181325953000876
1287 0.172744928875239
1394 0.169711397542423
1510 0.165168112863284
1635 0.165322206212235
1771 0.179486727055505
1918 0.177499902113277
2078 0.181573937061172
2250 0.182260098277007
2437 0.171506667247533
2640 0.168201046072913
2859 0.152023238575331
3097 0.155422356811239
3354 0.155193483787524
3633 0.154252249184797
3935 0.140909746785724
4262 0.141282622652018
4616 0.128953716819263
5000 0.130361483791334
};
\addplot [line width=0.7000000000000001pt, black, dashed]
table {%
100 0.481364381859579
108 0.464755356694629
117 0.457494316870976
127 0.442958436477485
138 0.426965188024608
149 0.414614453029318
161 0.407976328043523
175 0.39725704373531
189 0.382483201693791
205 0.385859304906856
222 0.369497595628203
241 0.355044028195492
261 0.3380847678872
282 0.325353057253603
306 0.321464161136803
331 0.32700287636065
359 0.316724687825919
389 0.331722993032317
421 0.323197185465412
456 0.322278622486992
494 0.312073396763021
535 0.317061922891473
579 0.308485945169615
627 0.306149332868807
679 0.295514190576037
736 0.294761943585706
797 0.289053331019986
863 0.288527190861965
935 0.277416827257005
1013 0.270282189571542
1097 0.270107424486116
1188 0.266672971052669
1287 0.251826057832938
1394 0.249748440634064
1510 0.241650768377245
1635 0.239740744554738
1771 0.24373690213204
1918 0.252909698876046
2078 0.247765724960661
2250 0.240223079383641
2437 0.232685578181925
2640 0.23580710707978
2859 0.238777227365712
3097 0.229057424148005
3354 0.22657521666947
3633 0.226081512867401
3935 0.226503357699914
4262 0.2221922417631
4616 0.218478444785539
5000 0.214155431920472
};
\end{axis}

\end{tikzpicture}%
			\subcaption{Bandwidth $h$.}
			\label{fig:bandwidth}
		\end{minipage}
		\caption{Left: The log-likelihood \cref{eq:loglikelihood estimation} with $m=10$ and the estimated entropy \cref{eq:entropy estimation} using the features $\Delta v$ and a bandwidth estimated using cross validation. Right: comparison of the bandwidth computed using the Silverman rule (for which the resulting log-likelihood is shown in \cref{fig:likelihood activity based}) and using cross validation (for which the result is shown in the left figure).}
		\label{fig:loglikelihood and bandwidth}
	\end{figure}
\end{example}


\subsection{Dissimilarity of new activities}
\label{sec:dissimilarity new activities}

This method compares the newly collected activities with the earlier collected activities. If the newly collected activities are very ``similar'' to the earlier collected activities, it is concluded that the dataset is complete. The main difference with the previous methods is that this method focuses on having distinct activities, while the distribution of the activities (i.e., the PDFs) is not taken into account. 

For determining the dissimilarity between two activities, the features of the activity are used. The feature vector is denoted as $\theta$. For this study, $\theta = \begin{bmatrix} v_{\textup{end}} & \Delta v & \Delta t \end{bmatrix}^T$, where $v_{\textup{end}}$, $\Delta v$, and $\Delta t$ are defined in \cref{sec:comparing pdf activities}.

A max-minimum method is used to determine to what degree the newly collected activities differ from the earlier collected activities, i.e., 
\begin{equation} \label{eq:maxmin dissimilarity}
	J_d(n, m) = \max_{n+1 \leq j \leq n+m} \min_{1 \leq i \leq n} g(\theta_i, \theta_j),
\end{equation}
where $g(\theta_i, \theta_j)$ denotes the dissimilarity measure that is defined as follows:
\begin{equation}
	g(\theta_i, \theta_j) = \sqrt{ \left( \theta_i - \theta_j \right)^T W \left( \theta_i - \theta_j \right)},
\end{equation}
where $W$ is a positive definite weighting matrix. For this study, $W$ is chosen to be a diagonal matrix with the $(k,k)$-th element equal to the reciprocal of the variance of the $k$-th feature, based on the $n$ activities.

\begin{example}
	Based on the artificial data, the max-minimum of the dissimilarity between the $m=10$ new activities and the earlier activities is computed, see \cref{fig:jd}. Although it is not clearly visible, there is a trend where $J_d(n,m)$ decreases as more data is added. There are, however, still some high values of $J_d(n,m)$ near the end. Furthermore, there is not a clear point in time where \emph{completeness} is reached.
	
	\setlength\figurewidth{0.6\linewidth}
	\setlength\figureheight{0.7\figurewidth}
	\begin{figure}
		\centering
		% This file was created by matplotlib2tikz v0.6.14.
\begin{tikzpicture}

\definecolor{color0}{rgb}{0.12156862745098,0.466666666666667,0.705882352941177}

\begin{axis}[
xlabel={Time [min]},
ylabel={$J_d(n,m)$},
xmin=-10.71025, xmax=295.901916666667,
ymin=-0.112508765448454, ymax=4.68639056662771,
width=\figurewidth,
height=\figureheight,
tick align=outside,
tick pos=left,
xmajorgrids,
x grid style={lightgray!92.026143790849673!black},
ymajorgrids,
y grid style={lightgray!92.026143790849673!black}
]
\addplot [semithick, color0, forget plot]
table {%
3.22666666666667 3.05527209954011
7.03166666666667 0.8952748275371
9.51333333333333 2.89002982683614
13.2833333333333 1.39555592070179
15.2183333333333 2.91873621800697
19.3116666666667 0.589163860592834
20.955 1.13597890170659
25.4433333333333 0.548546046571543
28.875 0.65341430886112
32.4783333333333 1.14977993632803
35.8016666666667 0.796981124717321
38.555 0.632204476294903
41.2883333333333 0.869887025551095
43.7 0.47888662974389
45.8583333333333 0.86387510649373
48.9416666666667 1.50816229656601
52.62 0.275572845326598
56.0516666666667 0.420177277288686
58.6816666666667 0.955071785457112
61.9083333333333 0.315024767740099
64.215 0.703759891700535
66.5166666666667 0.384155105235553
69.2616666666667 4.46825877880607
74.175 0.368473861251218
76.2166666666667 1.52913432196145
80.58 0.800613612437419
82.835 2.11910075840315
86.845 0.706027151844831
89.09 1.21223335399281
93.72 1.21287733783914
97.545 0.383505085750361
100.02 0.791230661884247
102.33 0.499084730199474
105.761666666667 0.296402645549264
109.115 0.162291639068505
112.37 0.313390017560296
115.588333333333 2.42811038156223
119.293333333333 0.470197764696669
121.851666666667 1.37437408053038
124.501666666667 0.976097737604506
127.805 0.291557581024465
130.513333333333 0.317816567490708
132.728333333333 0.226197274389098
136.228333333333 0.369076325952674
138.8 0.478245650845436
140.811666666667 0.166560106262237
142.7 0.195818343577712
145.056666666667 0.38510926736786
147.041666666667 0.212040519208788
149.461666666667 0.392050460349653
151.1 0.328483300087497
152.615 0.271196797525988
154.468333333333 0.187438607759445
156.353333333333 0.352039307491664
158.696666666667 0.209031733781496
160.356666666667 0.351174339403669
162.825 0.212936519315872
165.365 0.128167621276153
168.71 0.693136846234701
172.803333333333 0.10562302237319
176.4 0.344468177370442
178.83 0.691189880299499
182.02 0.17604920381747
184.616666666667 0.355750610887578
186.551666666667 0.816058591633282
189.38 0.55610727752017
192.69 0.421218853872525
195.335 0.355707170791158
198.151666666667 1.20039161123297
201.988333333333 0.927389855554418
205.976666666667 0.414940660157579
208.3 0.391969914736907
210.438333333333 0.520475903446518
215.353333333333 0.511162796432421
219.745 0.394789413327862
222.001666666667 0.678452082860962
225.611666666667 0.493684550050444
228.073333333333 0.141845861534957
230.448333333333 0.258594699283385
232.358333333333 1.61446896668767
236.365 0.648577206541811
239.905 0.39771984357716
242.72 0.321908233898065
247.211666666667 0.280814605433467
249.07 0.201634219087379
250.603333333333 0.334760342325922
252.301666666667 0.472049768746222
254.623333333333 0.522291430426178
256.77 0.81951613551443
259.576666666667 0.135591166475229
262.866666666667 0.200628836007872
265.441666666667 0.357639220337388
269.91 0.693007061467956
272.08 0.279362559199708
275.341666666667 0.710140481910438
277.97 0.300195518768077
279.988333333333 0.457564047448797
281.965 0.117582549490763
};
\end{axis}

\end{tikzpicture}%
		\caption{Dissimilarity \cref{eq:maxmin dissimilarity} of the $m=10$ new activities compared to the earlier recorded activities.}
		\label{fig:jd}
	\end{figure}
\end{example}

\printbibliography

\end{document} 